\documentclass[11pt]{article}

    \usepackage[breakable]{tcolorbox}
    \usepackage{parskip} % Stop auto-indenting (to mimic markdown behaviour)
    

    % Basic figure setup, for now with no caption control since it's done
    % automatically by Pandoc (which extracts ![](path) syntax from Markdown).
    \usepackage{graphicx}
    % Keep aspect ratio if custom image width or height is specified
    \setkeys{Gin}{keepaspectratio}
    % Maintain compatibility with old templates. Remove in nbconvert 6.0
    \let\Oldincludegraphics\includegraphics
    % Ensure that by default, figures have no caption (until we provide a
    % proper Figure object with a Caption API and a way to capture that
    % in the conversion process - todo).
    \usepackage{caption}
    \DeclareCaptionFormat{nocaption}{}
    \captionsetup{format=nocaption,aboveskip=0pt,belowskip=0pt}

    \usepackage{float}
    \floatplacement{figure}{H} % forces figures to be placed at the correct location
    \usepackage{xcolor} % Allow colors to be defined
    \usepackage{enumerate} % Needed for markdown enumerations to work
    \usepackage{geometry} % Used to adjust the document margins
    \usepackage{amsmath} % Equations
    \usepackage{amssymb} % Equations
    \usepackage{textcomp} % defines textquotesingle
    % Hack from http://tex.stackexchange.com/a/47451/13684:
    \AtBeginDocument{%
        \def\PYZsq{\textquotesingle}% Upright quotes in Pygmentized code
    }
    \usepackage{upquote} % Upright quotes for verbatim code
    \usepackage{eurosym} % defines \euro

    \usepackage{iftex}
    \ifPDFTeX
        \usepackage[T1]{fontenc}
        \IfFileExists{alphabeta.sty}{
              \usepackage{alphabeta}
          }{
              \usepackage[mathletters]{ucs}
              \usepackage[utf8x]{inputenc}
          }
    \else
        \usepackage{fontspec}
        \usepackage{unicode-math}
    \fi

    \usepackage{fancyvrb} % verbatim replacement that allows latex
    \usepackage{grffile} % extends the file name processing of package graphics
                         % to support a larger range
    \makeatletter % fix for old versions of grffile with XeLaTeX
    \@ifpackagelater{grffile}{2019/11/01}
    {
      % Do nothing on new versions
    }
    {
      \def\Gread@@xetex#1{%
        \IfFileExists{"\Gin@base".bb}%
        {\Gread@eps{\Gin@base.bb}}%
        {\Gread@@xetex@aux#1}%
      }
    }
    \makeatother
    \usepackage[Export]{adjustbox} % Used to constrain images to a maximum size
    \adjustboxset{max size={0.9\linewidth}{0.9\paperheight}}

    % The hyperref package gives us a pdf with properly built
    % internal navigation ('pdf bookmarks' for the table of contents,
    % internal cross-reference links, web links for URLs, etc.)
    \usepackage{hyperref}
    % The default LaTeX title has an obnoxious amount of whitespace. By default,
    % titling removes some of it. It also provides customization options.
    \usepackage{titling}
    \usepackage{longtable} % longtable support required by pandoc >1.10
    \usepackage{booktabs}  % table support for pandoc > 1.12.2
    \usepackage{array}     % table support for pandoc >= 2.11.3
    \usepackage{calc}      % table minipage width calculation for pandoc >= 2.11.1
    \usepackage[inline]{enumitem} % IRkernel/repr support (it uses the enumerate* environment)
    \usepackage[normalem]{ulem} % ulem is needed to support strikethroughs (\sout)
                                % normalem makes italics be italics, not underlines
    \usepackage{soul}      % strikethrough (\st) support for pandoc >= 3.0.0
    \usepackage{mathrsfs}
    

    
    % Colors for the hyperref package
    \definecolor{urlcolor}{rgb}{0,.145,.698}
    \definecolor{linkcolor}{rgb}{.71,0.21,0.01}
    \definecolor{citecolor}{rgb}{.12,.54,.11}

    % ANSI colors
    \definecolor{ansi-black}{HTML}{3E424D}
    \definecolor{ansi-black-intense}{HTML}{282C36}
    \definecolor{ansi-red}{HTML}{E75C58}
    \definecolor{ansi-red-intense}{HTML}{B22B31}
    \definecolor{ansi-green}{HTML}{00A250}
    \definecolor{ansi-green-intense}{HTML}{007427}
    \definecolor{ansi-yellow}{HTML}{DDB62B}
    \definecolor{ansi-yellow-intense}{HTML}{B27D12}
    \definecolor{ansi-blue}{HTML}{208FFB}
    \definecolor{ansi-blue-intense}{HTML}{0065CA}
    \definecolor{ansi-magenta}{HTML}{D160C4}
    \definecolor{ansi-magenta-intense}{HTML}{A03196}
    \definecolor{ansi-cyan}{HTML}{60C6C8}
    \definecolor{ansi-cyan-intense}{HTML}{258F8F}
    \definecolor{ansi-white}{HTML}{C5C1B4}
    \definecolor{ansi-white-intense}{HTML}{A1A6B2}
    \definecolor{ansi-default-inverse-fg}{HTML}{FFFFFF}
    \definecolor{ansi-default-inverse-bg}{HTML}{000000}

    % common color for the border for error outputs.
    \definecolor{outerrorbackground}{HTML}{FFDFDF}

    % commands and environments needed by pandoc snippets
    % extracted from the output of `pandoc -s`
    \providecommand{\tightlist}{%
      \setlength{\itemsep}{0pt}\setlength{\parskip}{0pt}}
    \DefineVerbatimEnvironment{Highlighting}{Verbatim}{commandchars=\\\{\}}
    % Add ',fontsize=\small' for more characters per line
    \newenvironment{Shaded}{}{}
    \newcommand{\KeywordTok}[1]{\textcolor[rgb]{0.00,0.44,0.13}{\textbf{{#1}}}}
    \newcommand{\DataTypeTok}[1]{\textcolor[rgb]{0.56,0.13,0.00}{{#1}}}
    \newcommand{\DecValTok}[1]{\textcolor[rgb]{0.25,0.63,0.44}{{#1}}}
    \newcommand{\BaseNTok}[1]{\textcolor[rgb]{0.25,0.63,0.44}{{#1}}}
    \newcommand{\FloatTok}[1]{\textcolor[rgb]{0.25,0.63,0.44}{{#1}}}
    \newcommand{\CharTok}[1]{\textcolor[rgb]{0.25,0.44,0.63}{{#1}}}
    \newcommand{\StringTok}[1]{\textcolor[rgb]{0.25,0.44,0.63}{{#1}}}
    \newcommand{\CommentTok}[1]{\textcolor[rgb]{0.38,0.63,0.69}{\textit{{#1}}}}
    \newcommand{\OtherTok}[1]{\textcolor[rgb]{0.00,0.44,0.13}{{#1}}}
    \newcommand{\AlertTok}[1]{\textcolor[rgb]{1.00,0.00,0.00}{\textbf{{#1}}}}
    \newcommand{\FunctionTok}[1]{\textcolor[rgb]{0.02,0.16,0.49}{{#1}}}
    \newcommand{\RegionMarkerTok}[1]{{#1}}
    \newcommand{\ErrorTok}[1]{\textcolor[rgb]{1.00,0.00,0.00}{\textbf{{#1}}}}
    \newcommand{\NormalTok}[1]{{#1}}

    % Additional commands for more recent versions of Pandoc
    \newcommand{\ConstantTok}[1]{\textcolor[rgb]{0.53,0.00,0.00}{{#1}}}
    \newcommand{\SpecialCharTok}[1]{\textcolor[rgb]{0.25,0.44,0.63}{{#1}}}
    \newcommand{\VerbatimStringTok}[1]{\textcolor[rgb]{0.25,0.44,0.63}{{#1}}}
    \newcommand{\SpecialStringTok}[1]{\textcolor[rgb]{0.73,0.40,0.53}{{#1}}}
    \newcommand{\ImportTok}[1]{{#1}}
    \newcommand{\DocumentationTok}[1]{\textcolor[rgb]{0.73,0.13,0.13}{\textit{{#1}}}}
    \newcommand{\AnnotationTok}[1]{\textcolor[rgb]{0.38,0.63,0.69}{\textbf{\textit{{#1}}}}}
    \newcommand{\CommentVarTok}[1]{\textcolor[rgb]{0.38,0.63,0.69}{\textbf{\textit{{#1}}}}}
    \newcommand{\VariableTok}[1]{\textcolor[rgb]{0.10,0.09,0.49}{{#1}}}
    \newcommand{\ControlFlowTok}[1]{\textcolor[rgb]{0.00,0.44,0.13}{\textbf{{#1}}}}
    \newcommand{\OperatorTok}[1]{\textcolor[rgb]{0.40,0.40,0.40}{{#1}}}
    \newcommand{\BuiltInTok}[1]{{#1}}
    \newcommand{\ExtensionTok}[1]{{#1}}
    \newcommand{\PreprocessorTok}[1]{\textcolor[rgb]{0.74,0.48,0.00}{{#1}}}
    \newcommand{\AttributeTok}[1]{\textcolor[rgb]{0.49,0.56,0.16}{{#1}}}
    \newcommand{\InformationTok}[1]{\textcolor[rgb]{0.38,0.63,0.69}{\textbf{\textit{{#1}}}}}
    \newcommand{\WarningTok}[1]{\textcolor[rgb]{0.38,0.63,0.69}{\textbf{\textit{{#1}}}}}
    \makeatletter
    \newsavebox\pandoc@box
    \newcommand*\pandocbounded[1]{%
      \sbox\pandoc@box{#1}%
      % scaling factors for width and height
      \Gscale@div\@tempa\textheight{\dimexpr\ht\pandoc@box+\dp\pandoc@box\relax}%
      \Gscale@div\@tempb\linewidth{\wd\pandoc@box}%
      % select the smaller of both
      \ifdim\@tempb\p@<\@tempa\p@
        \let\@tempa\@tempb
      \fi
      % scaling accordingly (\@tempa < 1)
      \ifdim\@tempa\p@<\p@
        \scalebox{\@tempa}{\usebox\pandoc@box}%
      % scaling not needed, use as it is
      \else
        \usebox{\pandoc@box}%
      \fi
    }
    \makeatother

    % Define a nice break command that doesn't care if a line doesn't already
    % exist.
    \def\br{\hspace*{\fill} \\* }
    % Math Jax compatibility definitions
    \def\gt{>}
    \def\lt{<}
    \let\Oldtex\TeX
    \let\Oldlatex\LaTeX
    \renewcommand{\TeX}{\textrm{\Oldtex}}
    \renewcommand{\LaTeX}{\textrm{\Oldlatex}}
    % Document parameters
    % Document title
    \title{Ch12-unsup-lab}
    
    
    
    
    
    
    
% Pygments definitions
\makeatletter
\def\PY@reset{\let\PY@it=\relax \let\PY@bf=\relax%
    \let\PY@ul=\relax \let\PY@tc=\relax%
    \let\PY@bc=\relax \let\PY@ff=\relax}
\def\PY@tok#1{\csname PY@tok@#1\endcsname}
\def\PY@toks#1+{\ifx\relax#1\empty\else%
    \PY@tok{#1}\expandafter\PY@toks\fi}
\def\PY@do#1{\PY@bc{\PY@tc{\PY@ul{%
    \PY@it{\PY@bf{\PY@ff{#1}}}}}}}
\def\PY#1#2{\PY@reset\PY@toks#1+\relax+\PY@do{#2}}

\@namedef{PY@tok@w}{\def\PY@tc##1{\textcolor[rgb]{0.73,0.73,0.73}{##1}}}
\@namedef{PY@tok@c}{\let\PY@it=\textit\def\PY@tc##1{\textcolor[rgb]{0.24,0.48,0.48}{##1}}}
\@namedef{PY@tok@cp}{\def\PY@tc##1{\textcolor[rgb]{0.61,0.40,0.00}{##1}}}
\@namedef{PY@tok@k}{\let\PY@bf=\textbf\def\PY@tc##1{\textcolor[rgb]{0.00,0.50,0.00}{##1}}}
\@namedef{PY@tok@kp}{\def\PY@tc##1{\textcolor[rgb]{0.00,0.50,0.00}{##1}}}
\@namedef{PY@tok@kt}{\def\PY@tc##1{\textcolor[rgb]{0.69,0.00,0.25}{##1}}}
\@namedef{PY@tok@o}{\def\PY@tc##1{\textcolor[rgb]{0.40,0.40,0.40}{##1}}}
\@namedef{PY@tok@ow}{\let\PY@bf=\textbf\def\PY@tc##1{\textcolor[rgb]{0.67,0.13,1.00}{##1}}}
\@namedef{PY@tok@nb}{\def\PY@tc##1{\textcolor[rgb]{0.00,0.50,0.00}{##1}}}
\@namedef{PY@tok@nf}{\def\PY@tc##1{\textcolor[rgb]{0.00,0.00,1.00}{##1}}}
\@namedef{PY@tok@nc}{\let\PY@bf=\textbf\def\PY@tc##1{\textcolor[rgb]{0.00,0.00,1.00}{##1}}}
\@namedef{PY@tok@nn}{\let\PY@bf=\textbf\def\PY@tc##1{\textcolor[rgb]{0.00,0.00,1.00}{##1}}}
\@namedef{PY@tok@ne}{\let\PY@bf=\textbf\def\PY@tc##1{\textcolor[rgb]{0.80,0.25,0.22}{##1}}}
\@namedef{PY@tok@nv}{\def\PY@tc##1{\textcolor[rgb]{0.10,0.09,0.49}{##1}}}
\@namedef{PY@tok@no}{\def\PY@tc##1{\textcolor[rgb]{0.53,0.00,0.00}{##1}}}
\@namedef{PY@tok@nl}{\def\PY@tc##1{\textcolor[rgb]{0.46,0.46,0.00}{##1}}}
\@namedef{PY@tok@ni}{\let\PY@bf=\textbf\def\PY@tc##1{\textcolor[rgb]{0.44,0.44,0.44}{##1}}}
\@namedef{PY@tok@na}{\def\PY@tc##1{\textcolor[rgb]{0.41,0.47,0.13}{##1}}}
\@namedef{PY@tok@nt}{\let\PY@bf=\textbf\def\PY@tc##1{\textcolor[rgb]{0.00,0.50,0.00}{##1}}}
\@namedef{PY@tok@nd}{\def\PY@tc##1{\textcolor[rgb]{0.67,0.13,1.00}{##1}}}
\@namedef{PY@tok@s}{\def\PY@tc##1{\textcolor[rgb]{0.73,0.13,0.13}{##1}}}
\@namedef{PY@tok@sd}{\let\PY@it=\textit\def\PY@tc##1{\textcolor[rgb]{0.73,0.13,0.13}{##1}}}
\@namedef{PY@tok@si}{\let\PY@bf=\textbf\def\PY@tc##1{\textcolor[rgb]{0.64,0.35,0.47}{##1}}}
\@namedef{PY@tok@se}{\let\PY@bf=\textbf\def\PY@tc##1{\textcolor[rgb]{0.67,0.36,0.12}{##1}}}
\@namedef{PY@tok@sr}{\def\PY@tc##1{\textcolor[rgb]{0.64,0.35,0.47}{##1}}}
\@namedef{PY@tok@ss}{\def\PY@tc##1{\textcolor[rgb]{0.10,0.09,0.49}{##1}}}
\@namedef{PY@tok@sx}{\def\PY@tc##1{\textcolor[rgb]{0.00,0.50,0.00}{##1}}}
\@namedef{PY@tok@m}{\def\PY@tc##1{\textcolor[rgb]{0.40,0.40,0.40}{##1}}}
\@namedef{PY@tok@gh}{\let\PY@bf=\textbf\def\PY@tc##1{\textcolor[rgb]{0.00,0.00,0.50}{##1}}}
\@namedef{PY@tok@gu}{\let\PY@bf=\textbf\def\PY@tc##1{\textcolor[rgb]{0.50,0.00,0.50}{##1}}}
\@namedef{PY@tok@gd}{\def\PY@tc##1{\textcolor[rgb]{0.63,0.00,0.00}{##1}}}
\@namedef{PY@tok@gi}{\def\PY@tc##1{\textcolor[rgb]{0.00,0.52,0.00}{##1}}}
\@namedef{PY@tok@gr}{\def\PY@tc##1{\textcolor[rgb]{0.89,0.00,0.00}{##1}}}
\@namedef{PY@tok@ge}{\let\PY@it=\textit}
\@namedef{PY@tok@gs}{\let\PY@bf=\textbf}
\@namedef{PY@tok@ges}{\let\PY@bf=\textbf\let\PY@it=\textit}
\@namedef{PY@tok@gp}{\let\PY@bf=\textbf\def\PY@tc##1{\textcolor[rgb]{0.00,0.00,0.50}{##1}}}
\@namedef{PY@tok@go}{\def\PY@tc##1{\textcolor[rgb]{0.44,0.44,0.44}{##1}}}
\@namedef{PY@tok@gt}{\def\PY@tc##1{\textcolor[rgb]{0.00,0.27,0.87}{##1}}}
\@namedef{PY@tok@err}{\def\PY@bc##1{{\setlength{\fboxsep}{\string -\fboxrule}\fcolorbox[rgb]{1.00,0.00,0.00}{1,1,1}{\strut ##1}}}}
\@namedef{PY@tok@kc}{\let\PY@bf=\textbf\def\PY@tc##1{\textcolor[rgb]{0.00,0.50,0.00}{##1}}}
\@namedef{PY@tok@kd}{\let\PY@bf=\textbf\def\PY@tc##1{\textcolor[rgb]{0.00,0.50,0.00}{##1}}}
\@namedef{PY@tok@kn}{\let\PY@bf=\textbf\def\PY@tc##1{\textcolor[rgb]{0.00,0.50,0.00}{##1}}}
\@namedef{PY@tok@kr}{\let\PY@bf=\textbf\def\PY@tc##1{\textcolor[rgb]{0.00,0.50,0.00}{##1}}}
\@namedef{PY@tok@bp}{\def\PY@tc##1{\textcolor[rgb]{0.00,0.50,0.00}{##1}}}
\@namedef{PY@tok@fm}{\def\PY@tc##1{\textcolor[rgb]{0.00,0.00,1.00}{##1}}}
\@namedef{PY@tok@vc}{\def\PY@tc##1{\textcolor[rgb]{0.10,0.09,0.49}{##1}}}
\@namedef{PY@tok@vg}{\def\PY@tc##1{\textcolor[rgb]{0.10,0.09,0.49}{##1}}}
\@namedef{PY@tok@vi}{\def\PY@tc##1{\textcolor[rgb]{0.10,0.09,0.49}{##1}}}
\@namedef{PY@tok@vm}{\def\PY@tc##1{\textcolor[rgb]{0.10,0.09,0.49}{##1}}}
\@namedef{PY@tok@sa}{\def\PY@tc##1{\textcolor[rgb]{0.73,0.13,0.13}{##1}}}
\@namedef{PY@tok@sb}{\def\PY@tc##1{\textcolor[rgb]{0.73,0.13,0.13}{##1}}}
\@namedef{PY@tok@sc}{\def\PY@tc##1{\textcolor[rgb]{0.73,0.13,0.13}{##1}}}
\@namedef{PY@tok@dl}{\def\PY@tc##1{\textcolor[rgb]{0.73,0.13,0.13}{##1}}}
\@namedef{PY@tok@s2}{\def\PY@tc##1{\textcolor[rgb]{0.73,0.13,0.13}{##1}}}
\@namedef{PY@tok@sh}{\def\PY@tc##1{\textcolor[rgb]{0.73,0.13,0.13}{##1}}}
\@namedef{PY@tok@s1}{\def\PY@tc##1{\textcolor[rgb]{0.73,0.13,0.13}{##1}}}
\@namedef{PY@tok@mb}{\def\PY@tc##1{\textcolor[rgb]{0.40,0.40,0.40}{##1}}}
\@namedef{PY@tok@mf}{\def\PY@tc##1{\textcolor[rgb]{0.40,0.40,0.40}{##1}}}
\@namedef{PY@tok@mh}{\def\PY@tc##1{\textcolor[rgb]{0.40,0.40,0.40}{##1}}}
\@namedef{PY@tok@mi}{\def\PY@tc##1{\textcolor[rgb]{0.40,0.40,0.40}{##1}}}
\@namedef{PY@tok@il}{\def\PY@tc##1{\textcolor[rgb]{0.40,0.40,0.40}{##1}}}
\@namedef{PY@tok@mo}{\def\PY@tc##1{\textcolor[rgb]{0.40,0.40,0.40}{##1}}}
\@namedef{PY@tok@ch}{\let\PY@it=\textit\def\PY@tc##1{\textcolor[rgb]{0.24,0.48,0.48}{##1}}}
\@namedef{PY@tok@cm}{\let\PY@it=\textit\def\PY@tc##1{\textcolor[rgb]{0.24,0.48,0.48}{##1}}}
\@namedef{PY@tok@cpf}{\let\PY@it=\textit\def\PY@tc##1{\textcolor[rgb]{0.24,0.48,0.48}{##1}}}
\@namedef{PY@tok@c1}{\let\PY@it=\textit\def\PY@tc##1{\textcolor[rgb]{0.24,0.48,0.48}{##1}}}
\@namedef{PY@tok@cs}{\let\PY@it=\textit\def\PY@tc##1{\textcolor[rgb]{0.24,0.48,0.48}{##1}}}

\def\PYZbs{\char`\\}
\def\PYZus{\char`\_}
\def\PYZob{\char`\{}
\def\PYZcb{\char`\}}
\def\PYZca{\char`\^}
\def\PYZam{\char`\&}
\def\PYZlt{\char`\<}
\def\PYZgt{\char`\>}
\def\PYZsh{\char`\#}
\def\PYZpc{\char`\%}
\def\PYZdl{\char`\$}
\def\PYZhy{\char`\-}
\def\PYZsq{\char`\'}
\def\PYZdq{\char`\"}
\def\PYZti{\char`\~}
% for compatibility with earlier versions
\def\PYZat{@}
\def\PYZlb{[}
\def\PYZrb{]}
\makeatother


    % For linebreaks inside Verbatim environment from package fancyvrb.
    \makeatletter
        \newbox\Wrappedcontinuationbox
        \newbox\Wrappedvisiblespacebox
        \newcommand*\Wrappedvisiblespace {\textcolor{red}{\textvisiblespace}}
        \newcommand*\Wrappedcontinuationsymbol {\textcolor{red}{\llap{\tiny$\m@th\hookrightarrow$}}}
        \newcommand*\Wrappedcontinuationindent {3ex }
        \newcommand*\Wrappedafterbreak {\kern\Wrappedcontinuationindent\copy\Wrappedcontinuationbox}
        % Take advantage of the already applied Pygments mark-up to insert
        % potential linebreaks for TeX processing.
        %        {, <, #, %, $, ' and ": go to next line.
        %        _, }, ^, &, >, - and ~: stay at end of broken line.
        % Use of \textquotesingle for straight quote.
        \newcommand*\Wrappedbreaksatspecials {%
            \def\PYGZus{\discretionary{\char`\_}{\Wrappedafterbreak}{\char`\_}}%
            \def\PYGZob{\discretionary{}{\Wrappedafterbreak\char`\{}{\char`\{}}%
            \def\PYGZcb{\discretionary{\char`\}}{\Wrappedafterbreak}{\char`\}}}%
            \def\PYGZca{\discretionary{\char`\^}{\Wrappedafterbreak}{\char`\^}}%
            \def\PYGZam{\discretionary{\char`\&}{\Wrappedafterbreak}{\char`\&}}%
            \def\PYGZlt{\discretionary{}{\Wrappedafterbreak\char`\<}{\char`\<}}%
            \def\PYGZgt{\discretionary{\char`\>}{\Wrappedafterbreak}{\char`\>}}%
            \def\PYGZsh{\discretionary{}{\Wrappedafterbreak\char`\#}{\char`\#}}%
            \def\PYGZpc{\discretionary{}{\Wrappedafterbreak\char`\%}{\char`\%}}%
            \def\PYGZdl{\discretionary{}{\Wrappedafterbreak\char`\$}{\char`\$}}%
            \def\PYGZhy{\discretionary{\char`\-}{\Wrappedafterbreak}{\char`\-}}%
            \def\PYGZsq{\discretionary{}{\Wrappedafterbreak\textquotesingle}{\textquotesingle}}%
            \def\PYGZdq{\discretionary{}{\Wrappedafterbreak\char`\"}{\char`\"}}%
            \def\PYGZti{\discretionary{\char`\~}{\Wrappedafterbreak}{\char`\~}}%
        }
        % Some characters . , ; ? ! / are not pygmentized.
        % This macro makes them "active" and they will insert potential linebreaks
        \newcommand*\Wrappedbreaksatpunct {%
            \lccode`\~`\.\lowercase{\def~}{\discretionary{\hbox{\char`\.}}{\Wrappedafterbreak}{\hbox{\char`\.}}}%
            \lccode`\~`\,\lowercase{\def~}{\discretionary{\hbox{\char`\,}}{\Wrappedafterbreak}{\hbox{\char`\,}}}%
            \lccode`\~`\;\lowercase{\def~}{\discretionary{\hbox{\char`\;}}{\Wrappedafterbreak}{\hbox{\char`\;}}}%
            \lccode`\~`\:\lowercase{\def~}{\discretionary{\hbox{\char`\:}}{\Wrappedafterbreak}{\hbox{\char`\:}}}%
            \lccode`\~`\?\lowercase{\def~}{\discretionary{\hbox{\char`\?}}{\Wrappedafterbreak}{\hbox{\char`\?}}}%
            \lccode`\~`\!\lowercase{\def~}{\discretionary{\hbox{\char`\!}}{\Wrappedafterbreak}{\hbox{\char`\!}}}%
            \lccode`\~`\/\lowercase{\def~}{\discretionary{\hbox{\char`\/}}{\Wrappedafterbreak}{\hbox{\char`\/}}}%
            \catcode`\.\active
            \catcode`\,\active
            \catcode`\;\active
            \catcode`\:\active
            \catcode`\?\active
            \catcode`\!\active
            \catcode`\/\active
            \lccode`\~`\~
        }
    \makeatother

    \let\OriginalVerbatim=\Verbatim
    \makeatletter
    \renewcommand{\Verbatim}[1][1]{%
        %\parskip\z@skip
        \sbox\Wrappedcontinuationbox {\Wrappedcontinuationsymbol}%
        \sbox\Wrappedvisiblespacebox {\FV@SetupFont\Wrappedvisiblespace}%
        \def\FancyVerbFormatLine ##1{\hsize\linewidth
            \vtop{\raggedright\hyphenpenalty\z@\exhyphenpenalty\z@
                \doublehyphendemerits\z@\finalhyphendemerits\z@
                \strut ##1\strut}%
        }%
        % If the linebreak is at a space, the latter will be displayed as visible
        % space at end of first line, and a continuation symbol starts next line.
        % Stretch/shrink are however usually zero for typewriter font.
        \def\FV@Space {%
            \nobreak\hskip\z@ plus\fontdimen3\font minus\fontdimen4\font
            \discretionary{\copy\Wrappedvisiblespacebox}{\Wrappedafterbreak}
            {\kern\fontdimen2\font}%
        }%

        % Allow breaks at special characters using \PYG... macros.
        \Wrappedbreaksatspecials
        % Breaks at punctuation characters . , ; ? ! and / need catcode=\active
        \OriginalVerbatim[#1,codes*=\Wrappedbreaksatpunct]%
    }
    \makeatother

    % Exact colors from NB
    \definecolor{incolor}{HTML}{303F9F}
    \definecolor{outcolor}{HTML}{D84315}
    \definecolor{cellborder}{HTML}{CFCFCF}
    \definecolor{cellbackground}{HTML}{F7F7F7}

    % prompt
    \makeatletter
    \newcommand{\boxspacing}{\kern\kvtcb@left@rule\kern\kvtcb@boxsep}
    \makeatother
    \newcommand{\prompt}[4]{
        {\ttfamily\llap{{\color{#2}[#3]:\hspace{3pt}#4}}\vspace{-\baselineskip}}
    }
    

    
    % Prevent overflowing lines due to hard-to-break entities
    \sloppy
    % Setup hyperref package
    \hypersetup{
      breaklinks=true,  % so long urls are correctly broken across lines
      colorlinks=true,
      urlcolor=urlcolor,
      linkcolor=linkcolor,
      citecolor=citecolor,
      }
    % Slightly bigger margins than the latex defaults
    
    \geometry{verbose,tmargin=1in,bmargin=1in,lmargin=1in,rmargin=1in}
    
    

\begin{document}
    
    \maketitle
    
    

    
    \hypertarget{unsupervised-learning}{%
\section{Unsupervised Learning}\label{unsupervised-learning}}

\href{https://mybinder.org/v2/gh/intro-stat-learning/ISLP_labs/v2.2?labpath=Ch12-unsup-lab.ipynb}{\includesvg{https://mybinder.org/badge_logo.svg}}

    In this lab we demonstrate PCA and clustering on several datasets. As in
other labs, we import some of our libraries at this top level. This
makes the code more readable, as scanning the first few lines of the
notebook tells us what libraries are used in this notebook.

    \begin{tcolorbox}[breakable, size=fbox, boxrule=1pt, pad at break*=1mm,colback=cellbackground, colframe=cellborder]
\prompt{In}{incolor}{1}{\boxspacing}
\begin{Verbatim}[commandchars=\\\{\}]
\PY{k+kn}{import}\PY{+w}{ }\PY{n+nn}{numpy}\PY{+w}{ }\PY{k}{as}\PY{+w}{ }\PY{n+nn}{np}
\PY{k+kn}{import}\PY{+w}{ }\PY{n+nn}{pandas}\PY{+w}{ }\PY{k}{as}\PY{+w}{ }\PY{n+nn}{pd}
\PY{k+kn}{import}\PY{+w}{ }\PY{n+nn}{matplotlib}\PY{n+nn}{.}\PY{n+nn}{pyplot}\PY{+w}{ }\PY{k}{as}\PY{+w}{ }\PY{n+nn}{plt}
\PY{k+kn}{from}\PY{+w}{ }\PY{n+nn}{statsmodels}\PY{n+nn}{.}\PY{n+nn}{datasets}\PY{+w}{ }\PY{k+kn}{import} \PY{n}{get\PYZus{}rdataset}
\PY{k+kn}{from}\PY{+w}{ }\PY{n+nn}{sklearn}\PY{n+nn}{.}\PY{n+nn}{decomposition}\PY{+w}{ }\PY{k+kn}{import} \PY{n}{PCA}
\PY{k+kn}{from}\PY{+w}{ }\PY{n+nn}{sklearn}\PY{n+nn}{.}\PY{n+nn}{preprocessing}\PY{+w}{ }\PY{k+kn}{import} \PY{n}{StandardScaler}
\PY{k+kn}{from}\PY{+w}{ }\PY{n+nn}{ISLP}\PY{+w}{ }\PY{k+kn}{import} \PY{n}{load\PYZus{}data}
\end{Verbatim}
\end{tcolorbox}

    We also collect the new imports needed for this lab.

    \begin{tcolorbox}[breakable, size=fbox, boxrule=1pt, pad at break*=1mm,colback=cellbackground, colframe=cellborder]
\prompt{In}{incolor}{2}{\boxspacing}
\begin{Verbatim}[commandchars=\\\{\}]
\PY{k+kn}{from}\PY{+w}{ }\PY{n+nn}{sklearn}\PY{n+nn}{.}\PY{n+nn}{cluster}\PY{+w}{ }\PY{k+kn}{import} \PYZbs{}
     \PY{p}{(}\PY{n}{KMeans}\PY{p}{,}
      \PY{n}{AgglomerativeClustering}\PY{p}{)}
\PY{k+kn}{from}\PY{+w}{ }\PY{n+nn}{scipy}\PY{n+nn}{.}\PY{n+nn}{cluster}\PY{n+nn}{.}\PY{n+nn}{hierarchy}\PY{+w}{ }\PY{k+kn}{import} \PYZbs{}
     \PY{p}{(}\PY{n}{dendrogram}\PY{p}{,}
      \PY{n}{cut\PYZus{}tree}\PY{p}{)}
\PY{k+kn}{from}\PY{+w}{ }\PY{n+nn}{ISLP}\PY{n+nn}{.}\PY{n+nn}{cluster}\PY{+w}{ }\PY{k+kn}{import} \PY{n}{compute\PYZus{}linkage}
\end{Verbatim}
\end{tcolorbox}

    \hypertarget{principal-components-analysis}{%
\subsection{Principal Components
Analysis}\label{principal-components-analysis}}

In this lab, we perform PCA on \texttt{USArrests}, a data set in the
\texttt{R} computing environment. We retrieve the data using
\texttt{get\_rdataset()}, which can fetch data from many standard
\texttt{R} packages.

The rows of the data set contain the 50 states, in alphabetical order.

    \begin{tcolorbox}[breakable, size=fbox, boxrule=1pt, pad at break*=1mm,colback=cellbackground, colframe=cellborder]
\prompt{In}{incolor}{3}{\boxspacing}
\begin{Verbatim}[commandchars=\\\{\}]
\PY{n}{USArrests} \PY{o}{=} \PY{n}{get\PYZus{}rdataset}\PY{p}{(}\PY{l+s+s1}{\PYZsq{}}\PY{l+s+s1}{USArrests}\PY{l+s+s1}{\PYZsq{}}\PY{p}{)}\PY{o}{.}\PY{n}{data}
\PY{n}{USArrests}
\end{Verbatim}
\end{tcolorbox}

            \begin{tcolorbox}[breakable, size=fbox, boxrule=.5pt, pad at break*=1mm, opacityfill=0]
\prompt{Out}{outcolor}{3}{\boxspacing}
\begin{Verbatim}[commandchars=\\\{\}]
                Murder  Assault  UrbanPop  Rape
rownames
Alabama           13.2      236        58  21.2
Alaska            10.0      263        48  44.5
Arizona            8.1      294        80  31.0
Arkansas           8.8      190        50  19.5
California         9.0      276        91  40.6
Colorado           7.9      204        78  38.7
Connecticut        3.3      110        77  11.1
Delaware           5.9      238        72  15.8
Florida           15.4      335        80  31.9
Georgia           17.4      211        60  25.8
Hawaii             5.3       46        83  20.2
Idaho              2.6      120        54  14.2
Illinois          10.4      249        83  24.0
Indiana            7.2      113        65  21.0
Iowa               2.2       56        57  11.3
Kansas             6.0      115        66  18.0
Kentucky           9.7      109        52  16.3
Louisiana         15.4      249        66  22.2
Maine              2.1       83        51   7.8
Maryland          11.3      300        67  27.8
Massachusetts      4.4      149        85  16.3
Michigan          12.1      255        74  35.1
Minnesota          2.7       72        66  14.9
Mississippi       16.1      259        44  17.1
Missouri           9.0      178        70  28.2
Montana            6.0      109        53  16.4
Nebraska           4.3      102        62  16.5
Nevada            12.2      252        81  46.0
New Hampshire      2.1       57        56   9.5
New Jersey         7.4      159        89  18.8
New Mexico        11.4      285        70  32.1
New York          11.1      254        86  26.1
North Carolina    13.0      337        45  16.1
North Dakota       0.8       45        44   7.3
Ohio               7.3      120        75  21.4
Oklahoma           6.6      151        68  20.0
Oregon             4.9      159        67  29.3
Pennsylvania       6.3      106        72  14.9
Rhode Island       3.4      174        87   8.3
South Carolina    14.4      279        48  22.5
South Dakota       3.8       86        45  12.8
Tennessee         13.2      188        59  26.9
Texas             12.7      201        80  25.5
Utah               3.2      120        80  22.9
Vermont            2.2       48        32  11.2
Virginia           8.5      156        63  20.7
Washington         4.0      145        73  26.2
West Virginia      5.7       81        39   9.3
Wisconsin          2.6       53        66  10.8
Wyoming            6.8      161        60  15.6
\end{Verbatim}
\end{tcolorbox}
        
    The columns of the data set contain the four variables.

    \begin{tcolorbox}[breakable, size=fbox, boxrule=1pt, pad at break*=1mm,colback=cellbackground, colframe=cellborder]
\prompt{In}{incolor}{4}{\boxspacing}
\begin{Verbatim}[commandchars=\\\{\}]
\PY{n}{USArrests}\PY{o}{.}\PY{n}{columns}
\end{Verbatim}
\end{tcolorbox}

            \begin{tcolorbox}[breakable, size=fbox, boxrule=.5pt, pad at break*=1mm, opacityfill=0]
\prompt{Out}{outcolor}{4}{\boxspacing}
\begin{Verbatim}[commandchars=\\\{\}]
Index(['Murder', 'Assault', 'UrbanPop', 'Rape'], dtype='object')
\end{Verbatim}
\end{tcolorbox}
        
    We first briefly examine the data. We notice that the variables have
vastly different means.

    \begin{tcolorbox}[breakable, size=fbox, boxrule=1pt, pad at break*=1mm,colback=cellbackground, colframe=cellborder]
\prompt{In}{incolor}{5}{\boxspacing}
\begin{Verbatim}[commandchars=\\\{\}]
\PY{n}{USArrests}\PY{o}{.}\PY{n}{mean}\PY{p}{(}\PY{p}{)}
\end{Verbatim}
\end{tcolorbox}

            \begin{tcolorbox}[breakable, size=fbox, boxrule=.5pt, pad at break*=1mm, opacityfill=0]
\prompt{Out}{outcolor}{5}{\boxspacing}
\begin{Verbatim}[commandchars=\\\{\}]
Murder        7.788
Assault     170.760
UrbanPop     65.540
Rape         21.232
dtype: float64
\end{Verbatim}
\end{tcolorbox}
        
    Dataframes have several useful methods for computing column-wise
summaries. We can also examine the variance of the four variables using
the \texttt{var()} method.

    \begin{tcolorbox}[breakable, size=fbox, boxrule=1pt, pad at break*=1mm,colback=cellbackground, colframe=cellborder]
\prompt{In}{incolor}{6}{\boxspacing}
\begin{Verbatim}[commandchars=\\\{\}]
\PY{n}{USArrests}\PY{o}{.}\PY{n}{var}\PY{p}{(}\PY{p}{)}
\end{Verbatim}
\end{tcolorbox}

            \begin{tcolorbox}[breakable, size=fbox, boxrule=.5pt, pad at break*=1mm, opacityfill=0]
\prompt{Out}{outcolor}{6}{\boxspacing}
\begin{Verbatim}[commandchars=\\\{\}]
Murder        18.970465
Assault     6945.165714
UrbanPop     209.518776
Rape          87.729159
dtype: float64
\end{Verbatim}
\end{tcolorbox}
        
    Not surprisingly, the variables also have vastly different variances.
The \texttt{UrbanPop} variable measures the percentage of the population
in each state living in an urban area, which is not a comparable number
to the number of rapes in each state per 100,000 individuals. PCA looks
for derived variables that account for most of the variance in the data
set. If we do not scale the variables before performing PCA, then the
principal components would mostly be driven by the \texttt{Assault}
variable, since it has by far the largest variance. So if the variables
are measured in different units or vary widely in scale, it is
recommended to standardize the variables to have standard deviation one
before performing PCA. Typically we set the means to zero as well.

This scaling can be done via the \texttt{StandardScaler()} transform
imported above. We first \texttt{fit} the scaler, which computes the
necessary means and standard deviations and then apply it to our data
using the \texttt{transform} method. As before, we combine these steps
using the \texttt{fit\_transform()} method.

    \begin{tcolorbox}[breakable, size=fbox, boxrule=1pt, pad at break*=1mm,colback=cellbackground, colframe=cellborder]
\prompt{In}{incolor}{7}{\boxspacing}
\begin{Verbatim}[commandchars=\\\{\}]
\PY{n}{scaler} \PY{o}{=} \PY{n}{StandardScaler}\PY{p}{(}\PY{n}{with\PYZus{}std}\PY{o}{=}\PY{k+kc}{True}\PY{p}{,}
                        \PY{n}{with\PYZus{}mean}\PY{o}{=}\PY{k+kc}{True}\PY{p}{)}
\PY{n}{USArrests\PYZus{}scaled} \PY{o}{=} \PY{n}{scaler}\PY{o}{.}\PY{n}{fit\PYZus{}transform}\PY{p}{(}\PY{n}{USArrests}\PY{p}{)}
\end{Verbatim}
\end{tcolorbox}

    Having scaled the data, we can then perform principal components
analysis using the \texttt{PCA()} transform from the
\texttt{sklearn.decomposition} package.

    \begin{tcolorbox}[breakable, size=fbox, boxrule=1pt, pad at break*=1mm,colback=cellbackground, colframe=cellborder]
\prompt{In}{incolor}{8}{\boxspacing}
\begin{Verbatim}[commandchars=\\\{\}]
\PY{n}{pcaUS} \PY{o}{=} \PY{n}{PCA}\PY{p}{(}\PY{p}{)}
\end{Verbatim}
\end{tcolorbox}

    (By default, the \texttt{PCA()} transform centers the variables to have
mean zero though it does not scale them.) The transform \texttt{pcaUS}
can be used to find the PCA \texttt{scores} returned by \texttt{fit()}.
Once the \texttt{fit} method has been called, the \texttt{pcaUS} object
also contains a number of useful quantities.

    \begin{tcolorbox}[breakable, size=fbox, boxrule=1pt, pad at break*=1mm,colback=cellbackground, colframe=cellborder]
\prompt{In}{incolor}{9}{\boxspacing}
\begin{Verbatim}[commandchars=\\\{\}]
\PY{n}{pcaUS}\PY{o}{.}\PY{n}{fit}\PY{p}{(}\PY{n}{USArrests\PYZus{}scaled}\PY{p}{)}
\end{Verbatim}
\end{tcolorbox}

            \begin{tcolorbox}[breakable, size=fbox, boxrule=.5pt, pad at break*=1mm, opacityfill=0]
\prompt{Out}{outcolor}{9}{\boxspacing}
\begin{Verbatim}[commandchars=\\\{\}]
PCA()
\end{Verbatim}
\end{tcolorbox}
        
    After fitting, the \texttt{mean\_} attribute corresponds to the means of
the variables. In this case, since we centered and scaled the data with
\texttt{scaler()} the means will all be 0.

    \begin{tcolorbox}[breakable, size=fbox, boxrule=1pt, pad at break*=1mm,colback=cellbackground, colframe=cellborder]
\prompt{In}{incolor}{10}{\boxspacing}
\begin{Verbatim}[commandchars=\\\{\}]
\PY{n}{pcaUS}\PY{o}{.}\PY{n}{mean\PYZus{}}
\end{Verbatim}
\end{tcolorbox}

            \begin{tcolorbox}[breakable, size=fbox, boxrule=.5pt, pad at break*=1mm, opacityfill=0]
\prompt{Out}{outcolor}{10}{\boxspacing}
\begin{Verbatim}[commandchars=\\\{\}]
array([-7.10542736e-17,  1.38777878e-16, -4.39648318e-16,  8.59312621e-16])
\end{Verbatim}
\end{tcolorbox}
        
    The scores can be computed using the \texttt{transform()} method of
\texttt{pcaUS} after it has been fit.

    \begin{tcolorbox}[breakable, size=fbox, boxrule=1pt, pad at break*=1mm,colback=cellbackground, colframe=cellborder]
\prompt{In}{incolor}{11}{\boxspacing}
\begin{Verbatim}[commandchars=\\\{\}]
\PY{n}{scores} \PY{o}{=} \PY{n}{pcaUS}\PY{o}{.}\PY{n}{transform}\PY{p}{(}\PY{n}{USArrests\PYZus{}scaled}\PY{p}{)}
\end{Verbatim}
\end{tcolorbox}

    We will plot these scores a bit further down. The \texttt{components\_}
attribute provides the principal component loadings: each row of
\texttt{pcaUS.components\_} contains the corresponding principal
component loading vector.

    \begin{tcolorbox}[breakable, size=fbox, boxrule=1pt, pad at break*=1mm,colback=cellbackground, colframe=cellborder]
\prompt{In}{incolor}{12}{\boxspacing}
\begin{Verbatim}[commandchars=\\\{\}]
\PY{n}{pcaUS}\PY{o}{.}\PY{n}{components\PYZus{}} 
\end{Verbatim}
\end{tcolorbox}

            \begin{tcolorbox}[breakable, size=fbox, boxrule=.5pt, pad at break*=1mm, opacityfill=0]
\prompt{Out}{outcolor}{12}{\boxspacing}
\begin{Verbatim}[commandchars=\\\{\}]
array([[ 0.53589947,  0.58318363,  0.27819087,  0.54343209],
       [-0.41818087, -0.1879856 ,  0.87280619,  0.16731864],
       [-0.34123273, -0.26814843, -0.37801579,  0.81777791],
       [-0.6492278 ,  0.74340748, -0.13387773, -0.08902432]])
\end{Verbatim}
\end{tcolorbox}
        
    The \texttt{biplot} is a common visualization method used with PCA. It
is not built in as a standard part of \texttt{sklearn}, though there are
python packages that do produce such plots. Here we make a simple biplot
manually.

    \begin{tcolorbox}[breakable, size=fbox, boxrule=1pt, pad at break*=1mm,colback=cellbackground, colframe=cellborder]
\prompt{In}{incolor}{13}{\boxspacing}
\begin{Verbatim}[commandchars=\\\{\}]
\PY{n}{i}\PY{p}{,} \PY{n}{j} \PY{o}{=} \PY{l+m+mi}{0}\PY{p}{,} \PY{l+m+mi}{1} \PY{c+c1}{\PYZsh{} which components}
\PY{n}{fig}\PY{p}{,} \PY{n}{ax} \PY{o}{=} \PY{n}{plt}\PY{o}{.}\PY{n}{subplots}\PY{p}{(}\PY{l+m+mi}{1}\PY{p}{,} \PY{l+m+mi}{1}\PY{p}{,} \PY{n}{figsize}\PY{o}{=}\PY{p}{(}\PY{l+m+mi}{8}\PY{p}{,} \PY{l+m+mi}{8}\PY{p}{)}\PY{p}{)}
\PY{n}{ax}\PY{o}{.}\PY{n}{scatter}\PY{p}{(}\PY{n}{scores}\PY{p}{[}\PY{p}{:}\PY{p}{,}\PY{l+m+mi}{0}\PY{p}{]}\PY{p}{,} \PY{n}{scores}\PY{p}{[}\PY{p}{:}\PY{p}{,}\PY{l+m+mi}{1}\PY{p}{]}\PY{p}{)}
\PY{n}{ax}\PY{o}{.}\PY{n}{set\PYZus{}xlabel}\PY{p}{(}\PY{l+s+s1}{\PYZsq{}}\PY{l+s+s1}{PC}\PY{l+s+si}{\PYZpc{}d}\PY{l+s+s1}{\PYZsq{}} \PY{o}{\PYZpc{}} \PY{p}{(}\PY{n}{i}\PY{o}{+}\PY{l+m+mi}{1}\PY{p}{)}\PY{p}{)}
\PY{n}{ax}\PY{o}{.}\PY{n}{set\PYZus{}ylabel}\PY{p}{(}\PY{l+s+s1}{\PYZsq{}}\PY{l+s+s1}{PC}\PY{l+s+si}{\PYZpc{}d}\PY{l+s+s1}{\PYZsq{}} \PY{o}{\PYZpc{}} \PY{p}{(}\PY{n}{j}\PY{o}{+}\PY{l+m+mi}{1}\PY{p}{)}\PY{p}{)}
\PY{k}{for} \PY{n}{k} \PY{o+ow}{in} \PY{n+nb}{range}\PY{p}{(}\PY{n}{pcaUS}\PY{o}{.}\PY{n}{components\PYZus{}}\PY{o}{.}\PY{n}{shape}\PY{p}{[}\PY{l+m+mi}{1}\PY{p}{]}\PY{p}{)}\PY{p}{:}
    \PY{n}{ax}\PY{o}{.}\PY{n}{arrow}\PY{p}{(}\PY{l+m+mi}{0}\PY{p}{,} \PY{l+m+mi}{0}\PY{p}{,} \PY{n}{pcaUS}\PY{o}{.}\PY{n}{components\PYZus{}}\PY{p}{[}\PY{n}{i}\PY{p}{,}\PY{n}{k}\PY{p}{]}\PY{p}{,} \PY{n}{pcaUS}\PY{o}{.}\PY{n}{components\PYZus{}}\PY{p}{[}\PY{n}{j}\PY{p}{,}\PY{n}{k}\PY{p}{]}\PY{p}{)}
    \PY{n}{ax}\PY{o}{.}\PY{n}{text}\PY{p}{(}\PY{n}{pcaUS}\PY{o}{.}\PY{n}{components\PYZus{}}\PY{p}{[}\PY{n}{i}\PY{p}{,}\PY{n}{k}\PY{p}{]}\PY{p}{,}
            \PY{n}{pcaUS}\PY{o}{.}\PY{n}{components\PYZus{}}\PY{p}{[}\PY{n}{j}\PY{p}{,}\PY{n}{k}\PY{p}{]}\PY{p}{,}
            \PY{n}{USArrests}\PY{o}{.}\PY{n}{columns}\PY{p}{[}\PY{n}{k}\PY{p}{]}\PY{p}{)}
\end{Verbatim}
\end{tcolorbox}

    \begin{center}
    \adjustimage{max size={0.9\linewidth}{0.9\paperheight}}{artifacts/latex/Ch12-unsup-lab_files/artifacts/latex/Ch12-unsup-lab_26_0.png}
    \end{center}
    { \hspace*{\fill} \\}
    
    Notice that this figure is a reflection of Figure 12.1 through the
\(y\)-axis. Recall that the principal components are only unique up to a
sign change, so we can reproduce that figure by flipping the signs of
the second set of scores and loadings. We also increase the length of
the arrows to emphasize the loadings.

    \begin{tcolorbox}[breakable, size=fbox, boxrule=1pt, pad at break*=1mm,colback=cellbackground, colframe=cellborder]
\prompt{In}{incolor}{14}{\boxspacing}
\begin{Verbatim}[commandchars=\\\{\}]
\PY{n}{scale\PYZus{}arrow} \PY{o}{=} \PY{n}{s\PYZus{}} \PY{o}{=} \PY{l+m+mi}{2}
\PY{n}{scores}\PY{p}{[}\PY{p}{:}\PY{p}{,}\PY{l+m+mi}{1}\PY{p}{]} \PY{o}{*}\PY{o}{=} \PY{o}{\PYZhy{}}\PY{l+m+mi}{1}
\PY{n}{pcaUS}\PY{o}{.}\PY{n}{components\PYZus{}}\PY{p}{[}\PY{l+m+mi}{1}\PY{p}{]} \PY{o}{*}\PY{o}{=} \PY{o}{\PYZhy{}}\PY{l+m+mi}{1} \PY{c+c1}{\PYZsh{} flip the y\PYZhy{}axis}
\PY{n}{fig}\PY{p}{,} \PY{n}{ax} \PY{o}{=} \PY{n}{plt}\PY{o}{.}\PY{n}{subplots}\PY{p}{(}\PY{l+m+mi}{1}\PY{p}{,} \PY{l+m+mi}{1}\PY{p}{,} \PY{n}{figsize}\PY{o}{=}\PY{p}{(}\PY{l+m+mi}{8}\PY{p}{,} \PY{l+m+mi}{8}\PY{p}{)}\PY{p}{)}
\PY{n}{ax}\PY{o}{.}\PY{n}{scatter}\PY{p}{(}\PY{n}{scores}\PY{p}{[}\PY{p}{:}\PY{p}{,}\PY{l+m+mi}{0}\PY{p}{]}\PY{p}{,} \PY{n}{scores}\PY{p}{[}\PY{p}{:}\PY{p}{,}\PY{l+m+mi}{1}\PY{p}{]}\PY{p}{)}
\PY{n}{ax}\PY{o}{.}\PY{n}{set\PYZus{}xlabel}\PY{p}{(}\PY{l+s+s1}{\PYZsq{}}\PY{l+s+s1}{PC}\PY{l+s+si}{\PYZpc{}d}\PY{l+s+s1}{\PYZsq{}} \PY{o}{\PYZpc{}} \PY{p}{(}\PY{n}{i}\PY{o}{+}\PY{l+m+mi}{1}\PY{p}{)}\PY{p}{)}
\PY{n}{ax}\PY{o}{.}\PY{n}{set\PYZus{}ylabel}\PY{p}{(}\PY{l+s+s1}{\PYZsq{}}\PY{l+s+s1}{PC}\PY{l+s+si}{\PYZpc{}d}\PY{l+s+s1}{\PYZsq{}} \PY{o}{\PYZpc{}} \PY{p}{(}\PY{n}{j}\PY{o}{+}\PY{l+m+mi}{1}\PY{p}{)}\PY{p}{)}
\PY{k}{for} \PY{n}{k} \PY{o+ow}{in} \PY{n+nb}{range}\PY{p}{(}\PY{n}{pcaUS}\PY{o}{.}\PY{n}{components\PYZus{}}\PY{o}{.}\PY{n}{shape}\PY{p}{[}\PY{l+m+mi}{1}\PY{p}{]}\PY{p}{)}\PY{p}{:}
    \PY{n}{ax}\PY{o}{.}\PY{n}{arrow}\PY{p}{(}\PY{l+m+mi}{0}\PY{p}{,} \PY{l+m+mi}{0}\PY{p}{,} \PY{n}{s\PYZus{}}\PY{o}{*}\PY{n}{pcaUS}\PY{o}{.}\PY{n}{components\PYZus{}}\PY{p}{[}\PY{n}{i}\PY{p}{,}\PY{n}{k}\PY{p}{]}\PY{p}{,} \PY{n}{s\PYZus{}}\PY{o}{*}\PY{n}{pcaUS}\PY{o}{.}\PY{n}{components\PYZus{}}\PY{p}{[}\PY{n}{j}\PY{p}{,}\PY{n}{k}\PY{p}{]}\PY{p}{)}
    \PY{n}{ax}\PY{o}{.}\PY{n}{text}\PY{p}{(}\PY{n}{s\PYZus{}}\PY{o}{*}\PY{n}{pcaUS}\PY{o}{.}\PY{n}{components\PYZus{}}\PY{p}{[}\PY{n}{i}\PY{p}{,}\PY{n}{k}\PY{p}{]}\PY{p}{,}
            \PY{n}{s\PYZus{}}\PY{o}{*}\PY{n}{pcaUS}\PY{o}{.}\PY{n}{components\PYZus{}}\PY{p}{[}\PY{n}{j}\PY{p}{,}\PY{n}{k}\PY{p}{]}\PY{p}{,}
            \PY{n}{USArrests}\PY{o}{.}\PY{n}{columns}\PY{p}{[}\PY{n}{k}\PY{p}{]}\PY{p}{)}
\end{Verbatim}
\end{tcolorbox}

    \begin{center}
    \adjustimage{max size={0.9\linewidth}{0.9\paperheight}}{artifacts/latex/Ch12-unsup-lab_files/artifacts/latex/Ch12-unsup-lab_28_0.png}
    \end{center}
    { \hspace*{\fill} \\}
    
    The standard deviations of the principal component scores are as
follows:

    \begin{tcolorbox}[breakable, size=fbox, boxrule=1pt, pad at break*=1mm,colback=cellbackground, colframe=cellborder]
\prompt{In}{incolor}{15}{\boxspacing}
\begin{Verbatim}[commandchars=\\\{\}]
\PY{n}{scores}\PY{o}{.}\PY{n}{std}\PY{p}{(}\PY{l+m+mi}{0}\PY{p}{,} \PY{n}{ddof}\PY{o}{=}\PY{l+m+mi}{1}\PY{p}{)}
\end{Verbatim}
\end{tcolorbox}

            \begin{tcolorbox}[breakable, size=fbox, boxrule=.5pt, pad at break*=1mm, opacityfill=0]
\prompt{Out}{outcolor}{15}{\boxspacing}
\begin{Verbatim}[commandchars=\\\{\}]
array([1.5908673 , 1.00496987, 0.6031915 , 0.4206774 ])
\end{Verbatim}
\end{tcolorbox}
        
    The variance of each score can be extracted directly from the
\texttt{pcaUS} object via the \texttt{explained\_variance\_} attribute.

    \begin{tcolorbox}[breakable, size=fbox, boxrule=1pt, pad at break*=1mm,colback=cellbackground, colframe=cellborder]
\prompt{In}{incolor}{16}{\boxspacing}
\begin{Verbatim}[commandchars=\\\{\}]
\PY{n}{pcaUS}\PY{o}{.}\PY{n}{explained\PYZus{}variance\PYZus{}}
\end{Verbatim}
\end{tcolorbox}

            \begin{tcolorbox}[breakable, size=fbox, boxrule=.5pt, pad at break*=1mm, opacityfill=0]
\prompt{Out}{outcolor}{16}{\boxspacing}
\begin{Verbatim}[commandchars=\\\{\}]
array([2.53085875, 1.00996444, 0.36383998, 0.17696948])
\end{Verbatim}
\end{tcolorbox}
        
    The proportion of variance explained by each principal component (PVE)
is stored as \texttt{explained\_variance\_ratio\_}:

    \begin{tcolorbox}[breakable, size=fbox, boxrule=1pt, pad at break*=1mm,colback=cellbackground, colframe=cellborder]
\prompt{In}{incolor}{17}{\boxspacing}
\begin{Verbatim}[commandchars=\\\{\}]
\PY{n}{pcaUS}\PY{o}{.}\PY{n}{explained\PYZus{}variance\PYZus{}ratio\PYZus{}}
\end{Verbatim}
\end{tcolorbox}

            \begin{tcolorbox}[breakable, size=fbox, boxrule=.5pt, pad at break*=1mm, opacityfill=0]
\prompt{Out}{outcolor}{17}{\boxspacing}
\begin{Verbatim}[commandchars=\\\{\}]
array([0.62006039, 0.24744129, 0.0891408 , 0.04335752])
\end{Verbatim}
\end{tcolorbox}
        
    We see that the first principal component explains 62.0\% of the
variance in the data, the next principal component explains 24.7\% of
the variance, and so forth. We can plot the PVE explained by each
component, as well as the cumulative PVE. We first plot the proportion
of variance explained.

    \begin{tcolorbox}[breakable, size=fbox, boxrule=1pt, pad at break*=1mm,colback=cellbackground, colframe=cellborder]
\prompt{In}{incolor}{18}{\boxspacing}
\begin{Verbatim}[commandchars=\\\{\}]
\PY{o}{\PYZpc{}\PYZpc{}capture}
\PY{n}{fig}\PY{p}{,} \PY{n}{axes} \PY{o}{=} \PY{n}{plt}\PY{o}{.}\PY{n}{subplots}\PY{p}{(}\PY{l+m+mi}{1}\PY{p}{,} \PY{l+m+mi}{2}\PY{p}{,} \PY{n}{figsize}\PY{o}{=}\PY{p}{(}\PY{l+m+mi}{15}\PY{p}{,} \PY{l+m+mi}{6}\PY{p}{)}\PY{p}{)}
\PY{n}{ticks} \PY{o}{=} \PY{n}{np}\PY{o}{.}\PY{n}{arange}\PY{p}{(}\PY{n}{pcaUS}\PY{o}{.}\PY{n}{n\PYZus{}components\PYZus{}}\PY{p}{)}\PY{o}{+}\PY{l+m+mi}{1}
\PY{n}{ax} \PY{o}{=} \PY{n}{axes}\PY{p}{[}\PY{l+m+mi}{0}\PY{p}{]}
\PY{n}{ax}\PY{o}{.}\PY{n}{plot}\PY{p}{(}\PY{n}{ticks}\PY{p}{,}
        \PY{n}{pcaUS}\PY{o}{.}\PY{n}{explained\PYZus{}variance\PYZus{}ratio\PYZus{}}\PY{p}{,}
        \PY{n}{marker}\PY{o}{=}\PY{l+s+s1}{\PYZsq{}}\PY{l+s+s1}{o}\PY{l+s+s1}{\PYZsq{}}\PY{p}{)}
\PY{n}{ax}\PY{o}{.}\PY{n}{set\PYZus{}xlabel}\PY{p}{(}\PY{l+s+s1}{\PYZsq{}}\PY{l+s+s1}{Principal Component}\PY{l+s+s1}{\PYZsq{}}\PY{p}{)}\PY{p}{;}
\PY{n}{ax}\PY{o}{.}\PY{n}{set\PYZus{}ylabel}\PY{p}{(}\PY{l+s+s1}{\PYZsq{}}\PY{l+s+s1}{Proportion of Variance Explained}\PY{l+s+s1}{\PYZsq{}}\PY{p}{)}
\PY{n}{ax}\PY{o}{.}\PY{n}{set\PYZus{}ylim}\PY{p}{(}\PY{p}{[}\PY{l+m+mi}{0}\PY{p}{,}\PY{l+m+mi}{1}\PY{p}{]}\PY{p}{)}
\PY{n}{ax}\PY{o}{.}\PY{n}{set\PYZus{}xticks}\PY{p}{(}\PY{n}{ticks}\PY{p}{)}
\end{Verbatim}
\end{tcolorbox}

    Notice the use of \texttt{\%\%capture}, which suppresses the displaying
of the partially completed figure.

    \begin{tcolorbox}[breakable, size=fbox, boxrule=1pt, pad at break*=1mm,colback=cellbackground, colframe=cellborder]
\prompt{In}{incolor}{19}{\boxspacing}
\begin{Verbatim}[commandchars=\\\{\}]
\PY{n}{ax} \PY{o}{=} \PY{n}{axes}\PY{p}{[}\PY{l+m+mi}{1}\PY{p}{]}
\PY{n}{ax}\PY{o}{.}\PY{n}{plot}\PY{p}{(}\PY{n}{ticks}\PY{p}{,}
        \PY{n}{pcaUS}\PY{o}{.}\PY{n}{explained\PYZus{}variance\PYZus{}ratio\PYZus{}}\PY{o}{.}\PY{n}{cumsum}\PY{p}{(}\PY{p}{)}\PY{p}{,}
        \PY{n}{marker}\PY{o}{=}\PY{l+s+s1}{\PYZsq{}}\PY{l+s+s1}{o}\PY{l+s+s1}{\PYZsq{}}\PY{p}{)}
\PY{n}{ax}\PY{o}{.}\PY{n}{set\PYZus{}xlabel}\PY{p}{(}\PY{l+s+s1}{\PYZsq{}}\PY{l+s+s1}{Principal Component}\PY{l+s+s1}{\PYZsq{}}\PY{p}{)}
\PY{n}{ax}\PY{o}{.}\PY{n}{set\PYZus{}ylabel}\PY{p}{(}\PY{l+s+s1}{\PYZsq{}}\PY{l+s+s1}{Cumulative Proportion of Variance Explained}\PY{l+s+s1}{\PYZsq{}}\PY{p}{)}
\PY{n}{ax}\PY{o}{.}\PY{n}{set\PYZus{}ylim}\PY{p}{(}\PY{p}{[}\PY{l+m+mi}{0}\PY{p}{,} \PY{l+m+mi}{1}\PY{p}{]}\PY{p}{)}
\PY{n}{ax}\PY{o}{.}\PY{n}{set\PYZus{}xticks}\PY{p}{(}\PY{n}{ticks}\PY{p}{)}
\PY{n}{fig}
\end{Verbatim}
\end{tcolorbox}
 
            
\prompt{Out}{outcolor}{19}{}
    
    \begin{center}
    \adjustimage{max size={0.9\linewidth}{0.9\paperheight}}{artifacts/latex/Ch12-unsup-lab_files/artifacts/latex/Ch12-unsup-lab_38_0.png}
    \end{center}
    { \hspace*{\fill} \\}
    

    The result is similar to that shown in Figure 12.3. Note that the method
\texttt{cumsum()} computes the cumulative sum of the elements of a
numeric vector. For instance:

    \begin{tcolorbox}[breakable, size=fbox, boxrule=1pt, pad at break*=1mm,colback=cellbackground, colframe=cellborder]
\prompt{In}{incolor}{20}{\boxspacing}
\begin{Verbatim}[commandchars=\\\{\}]
\PY{n}{a} \PY{o}{=} \PY{n}{np}\PY{o}{.}\PY{n}{array}\PY{p}{(}\PY{p}{[}\PY{l+m+mi}{1}\PY{p}{,}\PY{l+m+mi}{2}\PY{p}{,}\PY{l+m+mi}{8}\PY{p}{,}\PY{o}{\PYZhy{}}\PY{l+m+mi}{3}\PY{p}{]}\PY{p}{)}
\PY{n}{np}\PY{o}{.}\PY{n}{cumsum}\PY{p}{(}\PY{n}{a}\PY{p}{)}
\end{Verbatim}
\end{tcolorbox}

            \begin{tcolorbox}[breakable, size=fbox, boxrule=.5pt, pad at break*=1mm, opacityfill=0]
\prompt{Out}{outcolor}{20}{\boxspacing}
\begin{Verbatim}[commandchars=\\\{\}]
array([ 1,  3, 11,  8])
\end{Verbatim}
\end{tcolorbox}
        
    \hypertarget{matrix-completion}{%
\subsection{Matrix Completion}\label{matrix-completion}}

We now re-create the analysis carried out on the \texttt{USArrests} data
in Section 12.3.

We saw in Section 12.2.2 that solving the optimization
problem\textasciitilde(12.6) on a centered data matrix \(\bf X\) is
equivalent to computing the first \(M\) principal components of the
data. We use our scaled and centered \texttt{USArrests} data as
\(\bf X\) below. The \emph{singular value decomposition} (SVD) is a
general algorithm for solving (12.6).

    \begin{tcolorbox}[breakable, size=fbox, boxrule=1pt, pad at break*=1mm,colback=cellbackground, colframe=cellborder]
\prompt{In}{incolor}{21}{\boxspacing}
\begin{Verbatim}[commandchars=\\\{\}]
\PY{n}{X} \PY{o}{=} \PY{n}{USArrests\PYZus{}scaled}
\PY{n}{U}\PY{p}{,} \PY{n}{D}\PY{p}{,} \PY{n}{V} \PY{o}{=} \PY{n}{np}\PY{o}{.}\PY{n}{linalg}\PY{o}{.}\PY{n}{svd}\PY{p}{(}\PY{n}{X}\PY{p}{,} \PY{n}{full\PYZus{}matrices}\PY{o}{=}\PY{k+kc}{False}\PY{p}{)}
\PY{n}{U}\PY{o}{.}\PY{n}{shape}\PY{p}{,} \PY{n}{D}\PY{o}{.}\PY{n}{shape}\PY{p}{,} \PY{n}{V}\PY{o}{.}\PY{n}{shape}
\end{Verbatim}
\end{tcolorbox}

            \begin{tcolorbox}[breakable, size=fbox, boxrule=.5pt, pad at break*=1mm, opacityfill=0]
\prompt{Out}{outcolor}{21}{\boxspacing}
\begin{Verbatim}[commandchars=\\\{\}]
((50, 4), (4,), (4, 4))
\end{Verbatim}
\end{tcolorbox}
        
    The \texttt{np.linalg.svd()} function returns three components,
\texttt{U}, \texttt{D} and \texttt{V}. The matrix \texttt{V} is
equivalent to the loading matrix from principal components (up to an
unimportant sign flip). Using the \texttt{full\_matrices=False} option
ensures that for a tall matrix the shape of \texttt{U} is the same as
the shape of \texttt{X}.

    \begin{tcolorbox}[breakable, size=fbox, boxrule=1pt, pad at break*=1mm,colback=cellbackground, colframe=cellborder]
\prompt{In}{incolor}{22}{\boxspacing}
\begin{Verbatim}[commandchars=\\\{\}]
\PY{n}{V}
\end{Verbatim}
\end{tcolorbox}

            \begin{tcolorbox}[breakable, size=fbox, boxrule=.5pt, pad at break*=1mm, opacityfill=0]
\prompt{Out}{outcolor}{22}{\boxspacing}
\begin{Verbatim}[commandchars=\\\{\}]
array([[-0.53589947, -0.58318363, -0.27819087, -0.54343209],
       [-0.41818087, -0.1879856 ,  0.87280619,  0.16731864],
       [ 0.34123273,  0.26814843,  0.37801579, -0.81777791],
       [ 0.6492278 , -0.74340748,  0.13387773,  0.08902432]])
\end{Verbatim}
\end{tcolorbox}
        
    \begin{tcolorbox}[breakable, size=fbox, boxrule=1pt, pad at break*=1mm,colback=cellbackground, colframe=cellborder]
\prompt{In}{incolor}{23}{\boxspacing}
\begin{Verbatim}[commandchars=\\\{\}]
\PY{n}{pcaUS}\PY{o}{.}\PY{n}{components\PYZus{}}
\end{Verbatim}
\end{tcolorbox}

            \begin{tcolorbox}[breakable, size=fbox, boxrule=.5pt, pad at break*=1mm, opacityfill=0]
\prompt{Out}{outcolor}{23}{\boxspacing}
\begin{Verbatim}[commandchars=\\\{\}]
array([[ 0.53589947,  0.58318363,  0.27819087,  0.54343209],
       [ 0.41818087,  0.1879856 , -0.87280619, -0.16731864],
       [-0.34123273, -0.26814843, -0.37801579,  0.81777791],
       [-0.6492278 ,  0.74340748, -0.13387773, -0.08902432]])
\end{Verbatim}
\end{tcolorbox}
        
    The matrix \texttt{U} corresponds to a \emph{standardized} version of
the PCA score matrix (each column standardized to have sum-of-squares
one). If we multiply each column of \texttt{U} by the corresponding
element of \texttt{D}, we recover the PCA scores exactly (up to a
meaningless sign flip).

    \begin{tcolorbox}[breakable, size=fbox, boxrule=1pt, pad at break*=1mm,colback=cellbackground, colframe=cellborder]
\prompt{In}{incolor}{24}{\boxspacing}
\begin{Verbatim}[commandchars=\\\{\}]
\PY{p}{(}\PY{n}{U} \PY{o}{*} \PY{n}{D}\PY{p}{[}\PY{k+kc}{None}\PY{p}{,}\PY{p}{:}\PY{p}{]}\PY{p}{)}\PY{p}{[}\PY{p}{:}\PY{l+m+mi}{3}\PY{p}{]}
\end{Verbatim}
\end{tcolorbox}

            \begin{tcolorbox}[breakable, size=fbox, boxrule=.5pt, pad at break*=1mm, opacityfill=0]
\prompt{Out}{outcolor}{24}{\boxspacing}
\begin{Verbatim}[commandchars=\\\{\}]
array([[-0.98556588, -1.13339238,  0.44426879,  0.15626714],
       [-1.95013775, -1.07321326, -2.04000333, -0.43858344],
       [-1.76316354,  0.74595678, -0.05478082, -0.83465292]])
\end{Verbatim}
\end{tcolorbox}
        
    \begin{tcolorbox}[breakable, size=fbox, boxrule=1pt, pad at break*=1mm,colback=cellbackground, colframe=cellborder]
\prompt{In}{incolor}{25}{\boxspacing}
\begin{Verbatim}[commandchars=\\\{\}]
\PY{n}{scores}\PY{p}{[}\PY{p}{:}\PY{l+m+mi}{3}\PY{p}{]}
\end{Verbatim}
\end{tcolorbox}

            \begin{tcolorbox}[breakable, size=fbox, boxrule=.5pt, pad at break*=1mm, opacityfill=0]
\prompt{Out}{outcolor}{25}{\boxspacing}
\begin{Verbatim}[commandchars=\\\{\}]
array([[ 0.98556588,  1.13339238, -0.44426879, -0.15626714],
       [ 1.95013775,  1.07321326,  2.04000333,  0.43858344],
       [ 1.76316354, -0.74595678,  0.05478082,  0.83465292]])
\end{Verbatim}
\end{tcolorbox}
        
    While it would be possible to carry out this lab using the
\texttt{PCA()} estimator, here we use the \texttt{np.linalg.svd()}
function in order to illustrate its use.

We now omit 20 entries in the \(50\times 4\) data matrix at random. We
do so by first selecting 20 rows (states) at random, and then selecting
one of the four entries in each row at random. This ensures that every
row has at least three observed values.

    \begin{tcolorbox}[breakable, size=fbox, boxrule=1pt, pad at break*=1mm,colback=cellbackground, colframe=cellborder]
\prompt{In}{incolor}{26}{\boxspacing}
\begin{Verbatim}[commandchars=\\\{\}]
\PY{n}{n\PYZus{}omit} \PY{o}{=} \PY{l+m+mi}{20}
\PY{n}{np}\PY{o}{.}\PY{n}{random}\PY{o}{.}\PY{n}{seed}\PY{p}{(}\PY{l+m+mi}{15}\PY{p}{)}
\PY{n}{r\PYZus{}idx} \PY{o}{=} \PY{n}{np}\PY{o}{.}\PY{n}{random}\PY{o}{.}\PY{n}{choice}\PY{p}{(}\PY{n}{np}\PY{o}{.}\PY{n}{arange}\PY{p}{(}\PY{n}{X}\PY{o}{.}\PY{n}{shape}\PY{p}{[}\PY{l+m+mi}{0}\PY{p}{]}\PY{p}{)}\PY{p}{,}
                         \PY{n}{n\PYZus{}omit}\PY{p}{,}
                         \PY{n}{replace}\PY{o}{=}\PY{k+kc}{False}\PY{p}{)}
\PY{n}{c\PYZus{}idx} \PY{o}{=} \PY{n}{np}\PY{o}{.}\PY{n}{random}\PY{o}{.}\PY{n}{choice}\PY{p}{(}\PY{n}{np}\PY{o}{.}\PY{n}{arange}\PY{p}{(}\PY{n}{X}\PY{o}{.}\PY{n}{shape}\PY{p}{[}\PY{l+m+mi}{1}\PY{p}{]}\PY{p}{)}\PY{p}{,}
                         \PY{n}{n\PYZus{}omit}\PY{p}{,}
                         \PY{n}{replace}\PY{o}{=}\PY{k+kc}{True}\PY{p}{)}
\PY{n}{Xna} \PY{o}{=} \PY{n}{X}\PY{o}{.}\PY{n}{copy}\PY{p}{(}\PY{p}{)}
\PY{n}{Xna}\PY{p}{[}\PY{n}{r\PYZus{}idx}\PY{p}{,} \PY{n}{c\PYZus{}idx}\PY{p}{]} \PY{o}{=} \PY{n}{np}\PY{o}{.}\PY{n}{nan}
\end{Verbatim}
\end{tcolorbox}

    Here the array \texttt{r\_idx} contains 20 integers from 0 to 49; this
represents the states (rows of \texttt{X}) that are selected to contain
missing values. And \texttt{c\_idx} contains 20 integers from 0 to 3,
representing the features (columns in \texttt{X}) that contain the
missing values for each of the selected states.

We now write some code to implement Algorithm 12.1. We first write a
function that takes in a matrix, and returns an approximation to the
matrix using the \texttt{svd()} function. This will be needed in Step 2
of Algorithm 12.1.

    \begin{tcolorbox}[breakable, size=fbox, boxrule=1pt, pad at break*=1mm,colback=cellbackground, colframe=cellborder]
\prompt{In}{incolor}{27}{\boxspacing}
\begin{Verbatim}[commandchars=\\\{\}]
\PY{k}{def}\PY{+w}{ }\PY{n+nf}{low\PYZus{}rank}\PY{p}{(}\PY{n}{X}\PY{p}{,} \PY{n}{M}\PY{o}{=}\PY{l+m+mi}{1}\PY{p}{)}\PY{p}{:}
    \PY{n}{U}\PY{p}{,} \PY{n}{D}\PY{p}{,} \PY{n}{V} \PY{o}{=} \PY{n}{np}\PY{o}{.}\PY{n}{linalg}\PY{o}{.}\PY{n}{svd}\PY{p}{(}\PY{n}{X}\PY{p}{)}
    \PY{n}{L} \PY{o}{=} \PY{n}{U}\PY{p}{[}\PY{p}{:}\PY{p}{,}\PY{p}{:}\PY{n}{M}\PY{p}{]} \PY{o}{*} \PY{n}{D}\PY{p}{[}\PY{k+kc}{None}\PY{p}{,}\PY{p}{:}\PY{n}{M}\PY{p}{]}
    \PY{k}{return} \PY{n}{L}\PY{o}{.}\PY{n}{dot}\PY{p}{(}\PY{n}{V}\PY{p}{[}\PY{p}{:}\PY{n}{M}\PY{p}{]}\PY{p}{)}
\end{Verbatim}
\end{tcolorbox}

    To conduct Step 1 of the algorithm, we initialize \texttt{Xhat} --- this
is \(\tilde{\bf X}\) in Algorithm 12.1 --- by replacing the missing
values with the column means of the non-missing entries. These are
stored in \texttt{Xbar} below after running \texttt{np.nanmean()} over
the row axis. We make a copy so that when we assign values to
\texttt{Xhat} below we do not also overwrite the values in \texttt{Xna}.

    \begin{tcolorbox}[breakable, size=fbox, boxrule=1pt, pad at break*=1mm,colback=cellbackground, colframe=cellborder]
\prompt{In}{incolor}{28}{\boxspacing}
\begin{Verbatim}[commandchars=\\\{\}]
\PY{n}{Xhat} \PY{o}{=} \PY{n}{Xna}\PY{o}{.}\PY{n}{copy}\PY{p}{(}\PY{p}{)}
\PY{n}{Xbar} \PY{o}{=} \PY{n}{np}\PY{o}{.}\PY{n}{nanmean}\PY{p}{(}\PY{n}{Xhat}\PY{p}{,} \PY{n}{axis}\PY{o}{=}\PY{l+m+mi}{0}\PY{p}{)}
\PY{n}{Xhat}\PY{p}{[}\PY{n}{r\PYZus{}idx}\PY{p}{,} \PY{n}{c\PYZus{}idx}\PY{p}{]} \PY{o}{=} \PY{n}{Xbar}\PY{p}{[}\PY{n}{c\PYZus{}idx}\PY{p}{]}
\end{Verbatim}
\end{tcolorbox}

    Before we begin Step 2, we set ourselves up to measure the progress of
our iterations:

    \begin{tcolorbox}[breakable, size=fbox, boxrule=1pt, pad at break*=1mm,colback=cellbackground, colframe=cellborder]
\prompt{In}{incolor}{29}{\boxspacing}
\begin{Verbatim}[commandchars=\\\{\}]
\PY{n}{thresh} \PY{o}{=} \PY{l+m+mf}{1e\PYZhy{}7}
\PY{n}{rel\PYZus{}err} \PY{o}{=} \PY{l+m+mi}{1}
\PY{n}{count} \PY{o}{=} \PY{l+m+mi}{0}
\PY{n}{ismiss} \PY{o}{=} \PY{n}{np}\PY{o}{.}\PY{n}{isnan}\PY{p}{(}\PY{n}{Xna}\PY{p}{)}
\PY{n}{mssold} \PY{o}{=} \PY{n}{np}\PY{o}{.}\PY{n}{mean}\PY{p}{(}\PY{n}{Xhat}\PY{p}{[}\PY{o}{\PYZti{}}\PY{n}{ismiss}\PY{p}{]}\PY{o}{*}\PY{o}{*}\PY{l+m+mi}{2}\PY{p}{)}
\PY{n}{mss0} \PY{o}{=} \PY{n}{np}\PY{o}{.}\PY{n}{mean}\PY{p}{(}\PY{n}{Xna}\PY{p}{[}\PY{o}{\PYZti{}}\PY{n}{ismiss}\PY{p}{]}\PY{o}{*}\PY{o}{*}\PY{l+m+mi}{2}\PY{p}{)}
\end{Verbatim}
\end{tcolorbox}

    Here \texttt{ismiss} is a logical matrix with the same dimensions as
\texttt{Xna}; a given element is \texttt{True} if the corresponding
matrix element is missing. The notation \texttt{\textasciitilde{}ismiss}
negates this boolean vector. This is useful because it allows us to
access both the missing and non-missing entries. We store the mean of
the squared non-missing elements in \texttt{mss0}. We store the mean
squared error of the non-missing elements of the old version of
\texttt{Xhat} in \texttt{mssold} (which currently agrees with
\texttt{mss0}). We plan to store the mean squared error of the
non-missing elements of the current version of \texttt{Xhat} in
\texttt{mss}, and will then iterate Step 2 of Algorithm 12.1 until the
\emph{relative error}, defined as \texttt{(mssold\ -\ mss)\ /\ mss0},
falls below \texttt{thresh\ =\ 1e-7}. \{Algorithm 12.1 tells us to
iterate Step 2 until (12.14) is no longer decreasing. Determining
whether (12.14) is decreasing requires us only to keep track of
\texttt{mssold\ -\ mss}. However, in practice, we keep track of
\texttt{(mssold\ -\ mss)\ /\ mss0} instead: this makes it so that the
number of iterations required for Algorithm 12.1 to converge does not
depend on whether we multiplied the raw data \(\bf X\) by a constant
factor.\}

In Step 2(a) of Algorithm 12.1, we approximate \texttt{Xhat} using
\texttt{low\_rank()}; we call this \texttt{Xapp}. In Step 2(b), we use
\texttt{Xapp} to update the estimates for elements in \texttt{Xhat} that
are missing in \texttt{Xna}. Finally, in Step 2(c), we compute the
relative error. These three steps are contained in the following
\texttt{while} loop:

    \begin{tcolorbox}[breakable, size=fbox, boxrule=1pt, pad at break*=1mm,colback=cellbackground, colframe=cellborder]
\prompt{In}{incolor}{30}{\boxspacing}
\begin{Verbatim}[commandchars=\\\{\}]
\PY{k}{while} \PY{n}{rel\PYZus{}err} \PY{o}{\PYZgt{}} \PY{n}{thresh}\PY{p}{:}
    \PY{n}{count} \PY{o}{+}\PY{o}{=} \PY{l+m+mi}{1}
    \PY{c+c1}{\PYZsh{} Step 2(a)}
    \PY{n}{Xapp} \PY{o}{=} \PY{n}{low\PYZus{}rank}\PY{p}{(}\PY{n}{Xhat}\PY{p}{,} \PY{n}{M}\PY{o}{=}\PY{l+m+mi}{1}\PY{p}{)}
    \PY{c+c1}{\PYZsh{} Step 2(b)}
    \PY{n}{Xhat}\PY{p}{[}\PY{n}{ismiss}\PY{p}{]} \PY{o}{=} \PY{n}{Xapp}\PY{p}{[}\PY{n}{ismiss}\PY{p}{]}
    \PY{c+c1}{\PYZsh{} Step 2(c)}
    \PY{n}{mss} \PY{o}{=} \PY{n}{np}\PY{o}{.}\PY{n}{mean}\PY{p}{(}\PY{p}{(}\PY{p}{(}\PY{n}{Xna} \PY{o}{\PYZhy{}} \PY{n}{Xapp}\PY{p}{)}\PY{p}{[}\PY{o}{\PYZti{}}\PY{n}{ismiss}\PY{p}{]}\PY{p}{)}\PY{o}{*}\PY{o}{*}\PY{l+m+mi}{2}\PY{p}{)}
    \PY{n}{rel\PYZus{}err} \PY{o}{=} \PY{p}{(}\PY{n}{mssold} \PY{o}{\PYZhy{}} \PY{n}{mss}\PY{p}{)} \PY{o}{/} \PY{n}{mss0}
    \PY{n}{mssold} \PY{o}{=} \PY{n}{mss}
    \PY{n+nb}{print}\PY{p}{(}\PY{l+s+s2}{\PYZdq{}}\PY{l+s+s2}{Iteration: }\PY{l+s+si}{\PYZob{}0\PYZcb{}}\PY{l+s+s2}{, MSS:}\PY{l+s+si}{\PYZob{}1:.3f\PYZcb{}}\PY{l+s+s2}{, Rel.Err }\PY{l+s+si}{\PYZob{}2:.2e\PYZcb{}}\PY{l+s+s2}{\PYZdq{}}
          \PY{o}{.}\PY{n}{format}\PY{p}{(}\PY{n}{count}\PY{p}{,} \PY{n}{mss}\PY{p}{,} \PY{n}{rel\PYZus{}err}\PY{p}{)}\PY{p}{)}
\end{Verbatim}
\end{tcolorbox}

    \begin{Verbatim}[commandchars=\\\{\}]
Iteration: 1, MSS:0.395, Rel.Err 5.99e-01
Iteration: 2, MSS:0.382, Rel.Err 1.33e-02
Iteration: 3, MSS:0.381, Rel.Err 1.44e-03
Iteration: 4, MSS:0.381, Rel.Err 1.79e-04
Iteration: 5, MSS:0.381, Rel.Err 2.58e-05
Iteration: 6, MSS:0.381, Rel.Err 4.22e-06
Iteration: 7, MSS:0.381, Rel.Err 7.65e-07
Iteration: 8, MSS:0.381, Rel.Err 1.48e-07
Iteration: 9, MSS:0.381, Rel.Err 2.95e-08
    \end{Verbatim}

    We see that after eight iterations, the relative error has fallen below
\texttt{thresh\ =\ 1e-7}, and so the algorithm terminates. When this
happens, the mean squared error of the non-missing elements equals
0.381.

Finally, we compute the correlation between the 20 imputed values and
the actual values:

    \begin{tcolorbox}[breakable, size=fbox, boxrule=1pt, pad at break*=1mm,colback=cellbackground, colframe=cellborder]
\prompt{In}{incolor}{31}{\boxspacing}
\begin{Verbatim}[commandchars=\\\{\}]
\PY{n}{np}\PY{o}{.}\PY{n}{corrcoef}\PY{p}{(}\PY{n}{Xapp}\PY{p}{[}\PY{n}{ismiss}\PY{p}{]}\PY{p}{,} \PY{n}{X}\PY{p}{[}\PY{n}{ismiss}\PY{p}{]}\PY{p}{)}\PY{p}{[}\PY{l+m+mi}{0}\PY{p}{,}\PY{l+m+mi}{1}\PY{p}{]}
\end{Verbatim}
\end{tcolorbox}

            \begin{tcolorbox}[breakable, size=fbox, boxrule=.5pt, pad at break*=1mm, opacityfill=0]
\prompt{Out}{outcolor}{31}{\boxspacing}
\begin{Verbatim}[commandchars=\\\{\}]
0.711356743429736
\end{Verbatim}
\end{tcolorbox}
        
    In this lab, we implemented Algorithm 12.1 ourselves for didactic
purposes. However, a reader who wishes to apply matrix completion to
their data might look to more specialized \texttt{Python}
implementations.

    \hypertarget{clustering}{%
\subsection{Clustering}\label{clustering}}

    \hypertarget{k-means-clustering}{%
\subsubsection{\texorpdfstring{\(K\)-Means
Clustering}{K-Means Clustering}}\label{k-means-clustering}}

The estimator \texttt{sklearn.cluster.KMeans()} performs \(K\)-means
clustering in \texttt{Python}. We begin with a simple simulated example
in which there truly are two clusters in the data: the first 25
observations have a mean shift relative to the next 25 observations.

    \begin{tcolorbox}[breakable, size=fbox, boxrule=1pt, pad at break*=1mm,colback=cellbackground, colframe=cellborder]
\prompt{In}{incolor}{32}{\boxspacing}
\begin{Verbatim}[commandchars=\\\{\}]
\PY{n}{np}\PY{o}{.}\PY{n}{random}\PY{o}{.}\PY{n}{seed}\PY{p}{(}\PY{l+m+mi}{0}\PY{p}{)}\PY{p}{;}
\PY{n}{X} \PY{o}{=} \PY{n}{np}\PY{o}{.}\PY{n}{random}\PY{o}{.}\PY{n}{standard\PYZus{}normal}\PY{p}{(}\PY{p}{(}\PY{l+m+mi}{50}\PY{p}{,}\PY{l+m+mi}{2}\PY{p}{)}\PY{p}{)}\PY{p}{;}
\PY{n}{X}\PY{p}{[}\PY{p}{:}\PY{l+m+mi}{25}\PY{p}{,}\PY{l+m+mi}{0}\PY{p}{]} \PY{o}{+}\PY{o}{=} \PY{l+m+mi}{3}\PY{p}{;}
\PY{n}{X}\PY{p}{[}\PY{p}{:}\PY{l+m+mi}{25}\PY{p}{,}\PY{l+m+mi}{1}\PY{p}{]} \PY{o}{\PYZhy{}}\PY{o}{=} \PY{l+m+mi}{4}\PY{p}{;}
\end{Verbatim}
\end{tcolorbox}

    We now perform \(K\)-means clustering with \(K=2\).

    \begin{tcolorbox}[breakable, size=fbox, boxrule=1pt, pad at break*=1mm,colback=cellbackground, colframe=cellborder]
\prompt{In}{incolor}{33}{\boxspacing}
\begin{Verbatim}[commandchars=\\\{\}]
\PY{n}{kmeans} \PY{o}{=} \PY{n}{KMeans}\PY{p}{(}\PY{n}{n\PYZus{}clusters}\PY{o}{=}\PY{l+m+mi}{2}\PY{p}{,}
                \PY{n}{random\PYZus{}state}\PY{o}{=}\PY{l+m+mi}{2}\PY{p}{,}
                \PY{n}{n\PYZus{}init}\PY{o}{=}\PY{l+m+mi}{20}\PY{p}{)}\PY{o}{.}\PY{n}{fit}\PY{p}{(}\PY{n}{X}\PY{p}{)}
\end{Verbatim}
\end{tcolorbox}

    We specify \texttt{random\_state} to make the results reproducible. The
cluster assignments of the 50 observations are contained in
\texttt{kmeans.labels\_}.

    \begin{tcolorbox}[breakable, size=fbox, boxrule=1pt, pad at break*=1mm,colback=cellbackground, colframe=cellborder]
\prompt{In}{incolor}{34}{\boxspacing}
\begin{Verbatim}[commandchars=\\\{\}]
\PY{n}{kmeans}\PY{o}{.}\PY{n}{labels\PYZus{}}
\end{Verbatim}
\end{tcolorbox}

            \begin{tcolorbox}[breakable, size=fbox, boxrule=.5pt, pad at break*=1mm, opacityfill=0]
\prompt{Out}{outcolor}{34}{\boxspacing}
\begin{Verbatim}[commandchars=\\\{\}]
array([0, 0, 0, 0, 0, 0, 0, 0, 0, 0, 0, 0, 0, 0, 0, 0, 0, 0, 0, 0, 0, 1,
       0, 0, 0, 1, 1, 1, 1, 1, 1, 1, 1, 1, 1, 1, 1, 1, 1, 1, 1, 1, 1, 1,
       1, 1, 1, 1, 1, 1], dtype=int32)
\end{Verbatim}
\end{tcolorbox}
        
    The \(K\)-means clustering perfectly separated the observations into two
clusters even though we did not supply any group information to
\texttt{KMeans()}. We can plot the data, with each observation colored
according to its cluster assignment.

    \begin{tcolorbox}[breakable, size=fbox, boxrule=1pt, pad at break*=1mm,colback=cellbackground, colframe=cellborder]
\prompt{In}{incolor}{35}{\boxspacing}
\begin{Verbatim}[commandchars=\\\{\}]
\PY{n}{fig}\PY{p}{,} \PY{n}{ax} \PY{o}{=} \PY{n}{plt}\PY{o}{.}\PY{n}{subplots}\PY{p}{(}\PY{l+m+mi}{1}\PY{p}{,} \PY{l+m+mi}{1}\PY{p}{,} \PY{n}{figsize}\PY{o}{=}\PY{p}{(}\PY{l+m+mi}{8}\PY{p}{,}\PY{l+m+mi}{8}\PY{p}{)}\PY{p}{)}
\PY{n}{ax}\PY{o}{.}\PY{n}{scatter}\PY{p}{(}\PY{n}{X}\PY{p}{[}\PY{p}{:}\PY{p}{,}\PY{l+m+mi}{0}\PY{p}{]}\PY{p}{,} \PY{n}{X}\PY{p}{[}\PY{p}{:}\PY{p}{,}\PY{l+m+mi}{1}\PY{p}{]}\PY{p}{,} \PY{n}{c}\PY{o}{=}\PY{n}{kmeans}\PY{o}{.}\PY{n}{labels\PYZus{}}\PY{p}{)}
\PY{n}{ax}\PY{o}{.}\PY{n}{set\PYZus{}title}\PY{p}{(}\PY{l+s+s2}{\PYZdq{}}\PY{l+s+s2}{K\PYZhy{}Means Clustering Results with K=2}\PY{l+s+s2}{\PYZdq{}}\PY{p}{)}\PY{p}{;}
\end{Verbatim}
\end{tcolorbox}

    \begin{center}
    \adjustimage{max size={0.9\linewidth}{0.9\paperheight}}{artifacts/latex/Ch12-unsup-lab_files/artifacts/latex/Ch12-unsup-lab_70_0.png}
    \end{center}
    { \hspace*{\fill} \\}
    
    Here the observations can be easily plotted because they are
two-dimensional. If there were more than two variables then we could
instead perform PCA and plot the first two principal component score
vectors to represent the clusters.

In this example, we knew that there really were two clusters because we
generated the data. However, for real data, we do not know the true
number of clusters, nor whether they exist in any precise way. We could
instead have performed \(K\)-means clustering on this example with
\(K=3\).

    \begin{tcolorbox}[breakable, size=fbox, boxrule=1pt, pad at break*=1mm,colback=cellbackground, colframe=cellborder]
\prompt{In}{incolor}{36}{\boxspacing}
\begin{Verbatim}[commandchars=\\\{\}]
\PY{n}{kmeans} \PY{o}{=} \PY{n}{KMeans}\PY{p}{(}\PY{n}{n\PYZus{}clusters}\PY{o}{=}\PY{l+m+mi}{3}\PY{p}{,}
                \PY{n}{random\PYZus{}state}\PY{o}{=}\PY{l+m+mi}{3}\PY{p}{,}
                \PY{n}{n\PYZus{}init}\PY{o}{=}\PY{l+m+mi}{20}\PY{p}{)}\PY{o}{.}\PY{n}{fit}\PY{p}{(}\PY{n}{X}\PY{p}{)}
\PY{n}{fig}\PY{p}{,} \PY{n}{ax} \PY{o}{=} \PY{n}{plt}\PY{o}{.}\PY{n}{subplots}\PY{p}{(}\PY{n}{figsize}\PY{o}{=}\PY{p}{(}\PY{l+m+mi}{8}\PY{p}{,}\PY{l+m+mi}{8}\PY{p}{)}\PY{p}{)}
\PY{n}{ax}\PY{o}{.}\PY{n}{scatter}\PY{p}{(}\PY{n}{X}\PY{p}{[}\PY{p}{:}\PY{p}{,}\PY{l+m+mi}{0}\PY{p}{]}\PY{p}{,} \PY{n}{X}\PY{p}{[}\PY{p}{:}\PY{p}{,}\PY{l+m+mi}{1}\PY{p}{]}\PY{p}{,} \PY{n}{c}\PY{o}{=}\PY{n}{kmeans}\PY{o}{.}\PY{n}{labels\PYZus{}}\PY{p}{)}
\PY{n}{ax}\PY{o}{.}\PY{n}{set\PYZus{}title}\PY{p}{(}\PY{l+s+s2}{\PYZdq{}}\PY{l+s+s2}{K\PYZhy{}Means Clustering Results with K=3}\PY{l+s+s2}{\PYZdq{}}\PY{p}{)}\PY{p}{;}
\end{Verbatim}
\end{tcolorbox}

    \begin{center}
    \adjustimage{max size={0.9\linewidth}{0.9\paperheight}}{artifacts/latex/Ch12-unsup-lab_files/artifacts/latex/Ch12-unsup-lab_72_0.png}
    \end{center}
    { \hspace*{\fill} \\}
    
    When \(K=3\), \(K\)-means clustering splits up the two clusters. We have
used the \texttt{n\_init} argument to run the \(K\)-means with 20
initial cluster assignments (the default is 10). If a value of
\texttt{n\_init} greater than one is used, then \(K\)-means clustering
will be performed using multiple random assignments in Step 1 of
Algorithm 12.2, and the \texttt{KMeans()} function will report only the
best results. Here we compare using \texttt{n\_init=1} to
\texttt{n\_init=20}.

    \begin{tcolorbox}[breakable, size=fbox, boxrule=1pt, pad at break*=1mm,colback=cellbackground, colframe=cellborder]
\prompt{In}{incolor}{37}{\boxspacing}
\begin{Verbatim}[commandchars=\\\{\}]
\PY{n}{kmeans1} \PY{o}{=} \PY{n}{KMeans}\PY{p}{(}\PY{n}{n\PYZus{}clusters}\PY{o}{=}\PY{l+m+mi}{3}\PY{p}{,}
                \PY{n}{random\PYZus{}state}\PY{o}{=}\PY{l+m+mi}{3}\PY{p}{,}
                \PY{n}{n\PYZus{}init}\PY{o}{=}\PY{l+m+mi}{1}\PY{p}{)}\PY{o}{.}\PY{n}{fit}\PY{p}{(}\PY{n}{X}\PY{p}{)}
\PY{n}{kmeans20} \PY{o}{=} \PY{n}{KMeans}\PY{p}{(}\PY{n}{n\PYZus{}clusters}\PY{o}{=}\PY{l+m+mi}{3}\PY{p}{,}
                  \PY{n}{random\PYZus{}state}\PY{o}{=}\PY{l+m+mi}{3}\PY{p}{,}
                  \PY{n}{n\PYZus{}init}\PY{o}{=}\PY{l+m+mi}{20}\PY{p}{)}\PY{o}{.}\PY{n}{fit}\PY{p}{(}\PY{n}{X}\PY{p}{)}\PY{p}{;}
\PY{n}{kmeans1}\PY{o}{.}\PY{n}{inertia\PYZus{}}\PY{p}{,} \PY{n}{kmeans20}\PY{o}{.}\PY{n}{inertia\PYZus{}}
\end{Verbatim}
\end{tcolorbox}

            \begin{tcolorbox}[breakable, size=fbox, boxrule=.5pt, pad at break*=1mm, opacityfill=0]
\prompt{Out}{outcolor}{37}{\boxspacing}
\begin{Verbatim}[commandchars=\\\{\}]
(76.85131986999252, 75.06261242745384)
\end{Verbatim}
\end{tcolorbox}
        
    Note that \texttt{kmeans.inertia\_} is the total within-cluster sum of
squares, which we seek to minimize by performing \(K\)-means clustering
(12.17).

We \emph{strongly} recommend always running \(K\)-means clustering with
a large value of \texttt{n\_init}, such as 20 or 50, since otherwise an
undesirable local optimum may be obtained.

When performing \(K\)-means clustering, in addition to using multiple
initial cluster assignments, it is also important to set a random seed
using the \texttt{random\_state} argument to \texttt{KMeans()}. This
way, the initial cluster assignments in Step 1 can be replicated, and
the \(K\)-means output will be fully reproducible.

    \hypertarget{hierarchical-clustering}{%
\subsubsection{Hierarchical Clustering}\label{hierarchical-clustering}}

The \texttt{AgglomerativeClustering()} class from the
\texttt{sklearn.clustering} package implements hierarchical clustering.
As its name is long, we use the short hand \texttt{HClust} for
\emph{hierarchical clustering}. Note that this will not change the
return type when using this method, so instances will still be of class
\texttt{AgglomerativeClustering}. In the following example we use the
data from the previous lab to plot the hierarchical clustering
dendrogram using complete, single, and average linkage clustering with
Euclidean distance as the dissimilarity measure. We begin by clustering
observations using complete linkage.

    \begin{tcolorbox}[breakable, size=fbox, boxrule=1pt, pad at break*=1mm,colback=cellbackground, colframe=cellborder]
\prompt{In}{incolor}{38}{\boxspacing}
\begin{Verbatim}[commandchars=\\\{\}]
\PY{n}{HClust} \PY{o}{=} \PY{n}{AgglomerativeClustering}
\PY{n}{hc\PYZus{}comp} \PY{o}{=} \PY{n}{HClust}\PY{p}{(}\PY{n}{distance\PYZus{}threshold}\PY{o}{=}\PY{l+m+mi}{0}\PY{p}{,}
                 \PY{n}{n\PYZus{}clusters}\PY{o}{=}\PY{k+kc}{None}\PY{p}{,}
                 \PY{n}{linkage}\PY{o}{=}\PY{l+s+s1}{\PYZsq{}}\PY{l+s+s1}{complete}\PY{l+s+s1}{\PYZsq{}}\PY{p}{)}
\PY{n}{hc\PYZus{}comp}\PY{o}{.}\PY{n}{fit}\PY{p}{(}\PY{n}{X}\PY{p}{)}
\end{Verbatim}
\end{tcolorbox}

            \begin{tcolorbox}[breakable, size=fbox, boxrule=.5pt, pad at break*=1mm, opacityfill=0]
\prompt{Out}{outcolor}{38}{\boxspacing}
\begin{Verbatim}[commandchars=\\\{\}]
AgglomerativeClustering(distance\_threshold=0, linkage='complete',
                        n\_clusters=None)
\end{Verbatim}
\end{tcolorbox}
        
    This computes the entire dendrogram. We could just as easily perform
hierarchical clustering with average or single linkage instead:

    \begin{tcolorbox}[breakable, size=fbox, boxrule=1pt, pad at break*=1mm,colback=cellbackground, colframe=cellborder]
\prompt{In}{incolor}{39}{\boxspacing}
\begin{Verbatim}[commandchars=\\\{\}]
\PY{n}{hc\PYZus{}avg} \PY{o}{=} \PY{n}{HClust}\PY{p}{(}\PY{n}{distance\PYZus{}threshold}\PY{o}{=}\PY{l+m+mi}{0}\PY{p}{,}
                \PY{n}{n\PYZus{}clusters}\PY{o}{=}\PY{k+kc}{None}\PY{p}{,}
                \PY{n}{linkage}\PY{o}{=}\PY{l+s+s1}{\PYZsq{}}\PY{l+s+s1}{average}\PY{l+s+s1}{\PYZsq{}}\PY{p}{)}\PY{p}{;}
\PY{n}{hc\PYZus{}avg}\PY{o}{.}\PY{n}{fit}\PY{p}{(}\PY{n}{X}\PY{p}{)}
\PY{n}{hc\PYZus{}sing} \PY{o}{=} \PY{n}{HClust}\PY{p}{(}\PY{n}{distance\PYZus{}threshold}\PY{o}{=}\PY{l+m+mi}{0}\PY{p}{,}
                 \PY{n}{n\PYZus{}clusters}\PY{o}{=}\PY{k+kc}{None}\PY{p}{,}
                 \PY{n}{linkage}\PY{o}{=}\PY{l+s+s1}{\PYZsq{}}\PY{l+s+s1}{single}\PY{l+s+s1}{\PYZsq{}}\PY{p}{)}\PY{p}{;}
\PY{n}{hc\PYZus{}sing}\PY{o}{.}\PY{n}{fit}\PY{p}{(}\PY{n}{X}\PY{p}{)}\PY{p}{;}
\end{Verbatim}
\end{tcolorbox}

    To use a precomputed distance matrix, we provide an additional argument
\texttt{metric="precomputed"}. In the code below, the first four lines
computes the \(50\times 50\) pairwise-distance matrix.

    \begin{tcolorbox}[breakable, size=fbox, boxrule=1pt, pad at break*=1mm,colback=cellbackground, colframe=cellborder]
\prompt{In}{incolor}{40}{\boxspacing}
\begin{Verbatim}[commandchars=\\\{\}]
\PY{n}{D} \PY{o}{=} \PY{n}{np}\PY{o}{.}\PY{n}{zeros}\PY{p}{(}\PY{p}{(}\PY{n}{X}\PY{o}{.}\PY{n}{shape}\PY{p}{[}\PY{l+m+mi}{0}\PY{p}{]}\PY{p}{,} \PY{n}{X}\PY{o}{.}\PY{n}{shape}\PY{p}{[}\PY{l+m+mi}{0}\PY{p}{]}\PY{p}{)}\PY{p}{)}\PY{p}{;}
\PY{k}{for} \PY{n}{i} \PY{o+ow}{in} \PY{n+nb}{range}\PY{p}{(}\PY{n}{X}\PY{o}{.}\PY{n}{shape}\PY{p}{[}\PY{l+m+mi}{0}\PY{p}{]}\PY{p}{)}\PY{p}{:}
    \PY{n}{x\PYZus{}} \PY{o}{=} \PY{n}{np}\PY{o}{.}\PY{n}{multiply}\PY{o}{.}\PY{n}{outer}\PY{p}{(}\PY{n}{np}\PY{o}{.}\PY{n}{ones}\PY{p}{(}\PY{n}{X}\PY{o}{.}\PY{n}{shape}\PY{p}{[}\PY{l+m+mi}{0}\PY{p}{]}\PY{p}{)}\PY{p}{,} \PY{n}{X}\PY{p}{[}\PY{n}{i}\PY{p}{]}\PY{p}{)}
    \PY{n}{D}\PY{p}{[}\PY{n}{i}\PY{p}{]} \PY{o}{=} \PY{n}{np}\PY{o}{.}\PY{n}{sqrt}\PY{p}{(}\PY{n}{np}\PY{o}{.}\PY{n}{sum}\PY{p}{(}\PY{p}{(}\PY{n}{X} \PY{o}{\PYZhy{}} \PY{n}{x\PYZus{}}\PY{p}{)}\PY{o}{*}\PY{o}{*}\PY{l+m+mi}{2}\PY{p}{,} \PY{l+m+mi}{1}\PY{p}{)}\PY{p}{)}\PY{p}{;}
\PY{n}{hc\PYZus{}sing\PYZus{}pre} \PY{o}{=} \PY{n}{HClust}\PY{p}{(}\PY{n}{distance\PYZus{}threshold}\PY{o}{=}\PY{l+m+mi}{0}\PY{p}{,}
                     \PY{n}{n\PYZus{}clusters}\PY{o}{=}\PY{k+kc}{None}\PY{p}{,}
                     \PY{n}{metric}\PY{o}{=}\PY{l+s+s1}{\PYZsq{}}\PY{l+s+s1}{precomputed}\PY{l+s+s1}{\PYZsq{}}\PY{p}{,}
                     \PY{n}{linkage}\PY{o}{=}\PY{l+s+s1}{\PYZsq{}}\PY{l+s+s1}{single}\PY{l+s+s1}{\PYZsq{}}\PY{p}{)}
\PY{n}{hc\PYZus{}sing\PYZus{}pre}\PY{o}{.}\PY{n}{fit}\PY{p}{(}\PY{n}{D}\PY{p}{)}
\end{Verbatim}
\end{tcolorbox}

            \begin{tcolorbox}[breakable, size=fbox, boxrule=.5pt, pad at break*=1mm, opacityfill=0]
\prompt{Out}{outcolor}{40}{\boxspacing}
\begin{Verbatim}[commandchars=\\\{\}]
AgglomerativeClustering(distance\_threshold=0, linkage='single',
                        metric='precomputed', n\_clusters=None)
\end{Verbatim}
\end{tcolorbox}
        
    We use \texttt{dendrogram()} from \texttt{scipy.cluster.hierarchy} to
plot the dendrogram. However, \texttt{dendrogram()} expects a so-called
\emph{linkage-matrix representation} of the clustering, which is not
provided by \texttt{AgglomerativeClustering()}, but can be computed. The
function \texttt{compute\_linkage()} in the \texttt{ISLP.cluster}
package is provided for this purpose.

We can now plot the dendrograms. The numbers at the bottom of the plot
identify each observation. The \texttt{dendrogram()} function has a
default method to color different branches of the tree that suggests a
pre-defined cut of the tree at a particular depth. We prefer to
overwrite this default by setting this threshold to be infinite. Since
we want this behavior for many dendrograms, we store these values in a
dictionary \texttt{cargs} and pass this as keyword arguments using the
notation \texttt{**cargs}.

    \begin{tcolorbox}[breakable, size=fbox, boxrule=1pt, pad at break*=1mm,colback=cellbackground, colframe=cellborder]
\prompt{In}{incolor}{41}{\boxspacing}
\begin{Verbatim}[commandchars=\\\{\}]
\PY{n}{cargs} \PY{o}{=} \PY{p}{\PYZob{}}\PY{l+s+s1}{\PYZsq{}}\PY{l+s+s1}{color\PYZus{}threshold}\PY{l+s+s1}{\PYZsq{}}\PY{p}{:}\PY{o}{\PYZhy{}}\PY{n}{np}\PY{o}{.}\PY{n}{inf}\PY{p}{,}
         \PY{l+s+s1}{\PYZsq{}}\PY{l+s+s1}{above\PYZus{}threshold\PYZus{}color}\PY{l+s+s1}{\PYZsq{}}\PY{p}{:}\PY{l+s+s1}{\PYZsq{}}\PY{l+s+s1}{black}\PY{l+s+s1}{\PYZsq{}}\PY{p}{\PYZcb{}}
\PY{n}{linkage\PYZus{}comp} \PY{o}{=} \PY{n}{compute\PYZus{}linkage}\PY{p}{(}\PY{n}{hc\PYZus{}comp}\PY{p}{)}
\PY{n}{fig}\PY{p}{,} \PY{n}{ax} \PY{o}{=} \PY{n}{plt}\PY{o}{.}\PY{n}{subplots}\PY{p}{(}\PY{l+m+mi}{1}\PY{p}{,} \PY{l+m+mi}{1}\PY{p}{,} \PY{n}{figsize}\PY{o}{=}\PY{p}{(}\PY{l+m+mi}{8}\PY{p}{,} \PY{l+m+mi}{8}\PY{p}{)}\PY{p}{)}
\PY{n}{dendrogram}\PY{p}{(}\PY{n}{linkage\PYZus{}comp}\PY{p}{,}
           \PY{n}{ax}\PY{o}{=}\PY{n}{ax}\PY{p}{,}
           \PY{o}{*}\PY{o}{*}\PY{n}{cargs}\PY{p}{)}\PY{p}{;}
\end{Verbatim}
\end{tcolorbox}

    \begin{center}
    \adjustimage{max size={0.9\linewidth}{0.9\paperheight}}{artifacts/latex/Ch12-unsup-lab_files/artifacts/latex/Ch12-unsup-lab_83_0.png}
    \end{center}
    { \hspace*{\fill} \\}
    
    We may want to color branches of the tree above and below a
cut-threshold differently. This can be achieved by changing the
\texttt{color\_threshold}. Let's cut the tree at a height of 4, coloring
links that merge above 4 in black.

    \begin{tcolorbox}[breakable, size=fbox, boxrule=1pt, pad at break*=1mm,colback=cellbackground, colframe=cellborder]
\prompt{In}{incolor}{42}{\boxspacing}
\begin{Verbatim}[commandchars=\\\{\}]
\PY{n}{fig}\PY{p}{,} \PY{n}{ax} \PY{o}{=} \PY{n}{plt}\PY{o}{.}\PY{n}{subplots}\PY{p}{(}\PY{l+m+mi}{1}\PY{p}{,} \PY{l+m+mi}{1}\PY{p}{,} \PY{n}{figsize}\PY{o}{=}\PY{p}{(}\PY{l+m+mi}{8}\PY{p}{,} \PY{l+m+mi}{8}\PY{p}{)}\PY{p}{)}
\PY{n}{dendrogram}\PY{p}{(}\PY{n}{linkage\PYZus{}comp}\PY{p}{,}
           \PY{n}{ax}\PY{o}{=}\PY{n}{ax}\PY{p}{,}
           \PY{n}{color\PYZus{}threshold}\PY{o}{=}\PY{l+m+mi}{4}\PY{p}{,}
           \PY{n}{above\PYZus{}threshold\PYZus{}color}\PY{o}{=}\PY{l+s+s1}{\PYZsq{}}\PY{l+s+s1}{black}\PY{l+s+s1}{\PYZsq{}}\PY{p}{)}\PY{p}{;}
\end{Verbatim}
\end{tcolorbox}

    \begin{center}
    \adjustimage{max size={0.9\linewidth}{0.9\paperheight}}{artifacts/latex/Ch12-unsup-lab_files/artifacts/latex/Ch12-unsup-lab_85_0.png}
    \end{center}
    { \hspace*{\fill} \\}
    
    To determine the cluster labels for each observation associated with a
given cut of the dendrogram, we can use the \texttt{cut\_tree()}
function from \texttt{scipy.cluster.hierarchy}:

    \begin{tcolorbox}[breakable, size=fbox, boxrule=1pt, pad at break*=1mm,colback=cellbackground, colframe=cellborder]
\prompt{In}{incolor}{43}{\boxspacing}
\begin{Verbatim}[commandchars=\\\{\}]
\PY{n}{cut\PYZus{}tree}\PY{p}{(}\PY{n}{linkage\PYZus{}comp}\PY{p}{,} \PY{n}{n\PYZus{}clusters}\PY{o}{=}\PY{l+m+mi}{4}\PY{p}{)}\PY{o}{.}\PY{n}{T}
\end{Verbatim}
\end{tcolorbox}

            \begin{tcolorbox}[breakable, size=fbox, boxrule=.5pt, pad at break*=1mm, opacityfill=0]
\prompt{Out}{outcolor}{43}{\boxspacing}
\begin{Verbatim}[commandchars=\\\{\}]
array([[0, 1, 0, 0, 1, 1, 0, 1, 0, 0, 2, 0, 0, 0, 1, 1, 0, 0, 1, 0, 0, 2,
        0, 2, 2, 3, 2, 3, 3, 3, 3, 2, 3, 3, 3, 3, 2, 3, 3, 3, 3, 2, 3, 3,
        3, 3, 3, 3, 3, 3]])
\end{Verbatim}
\end{tcolorbox}
        
    This can also be achieved by providing an argument \texttt{n\_clusters}
to \texttt{HClust()}; however each cut would require recomputing the
clustering. Similarly, trees may be cut by distance threshold with an
argument of \texttt{distance\_threshold} to \texttt{HClust()} or
\texttt{height} to \texttt{cut\_tree()}.

    \begin{tcolorbox}[breakable, size=fbox, boxrule=1pt, pad at break*=1mm,colback=cellbackground, colframe=cellborder]
\prompt{In}{incolor}{44}{\boxspacing}
\begin{Verbatim}[commandchars=\\\{\}]
\PY{n}{cut\PYZus{}tree}\PY{p}{(}\PY{n}{linkage\PYZus{}comp}\PY{p}{,} \PY{n}{height}\PY{o}{=}\PY{l+m+mi}{5}\PY{p}{)}
\end{Verbatim}
\end{tcolorbox}

            \begin{tcolorbox}[breakable, size=fbox, boxrule=.5pt, pad at break*=1mm, opacityfill=0]
\prompt{Out}{outcolor}{44}{\boxspacing}
\begin{Verbatim}[commandchars=\\\{\}]
array([[0],
       [0],
       [0],
       [0],
       [0],
       [0],
       [0],
       [0],
       [0],
       [0],
       [1],
       [0],
       [0],
       [0],
       [0],
       [0],
       [0],
       [0],
       [0],
       [0],
       [0],
       [1],
       [0],
       [1],
       [1],
       [2],
       [1],
       [2],
       [2],
       [2],
       [2],
       [1],
       [2],
       [2],
       [2],
       [2],
       [1],
       [2],
       [2],
       [2],
       [2],
       [1],
       [2],
       [2],
       [2],
       [2],
       [2],
       [2],
       [2],
       [2]])
\end{Verbatim}
\end{tcolorbox}
        
    To scale the variables before performing hierarchical clustering of the
observations, we use \texttt{StandardScaler()} as in our PCA example:

    \begin{tcolorbox}[breakable, size=fbox, boxrule=1pt, pad at break*=1mm,colback=cellbackground, colframe=cellborder]
\prompt{In}{incolor}{45}{\boxspacing}
\begin{Verbatim}[commandchars=\\\{\}]
\PY{n}{scaler} \PY{o}{=} \PY{n}{StandardScaler}\PY{p}{(}\PY{p}{)}
\PY{n}{X\PYZus{}scale} \PY{o}{=} \PY{n}{scaler}\PY{o}{.}\PY{n}{fit\PYZus{}transform}\PY{p}{(}\PY{n}{X}\PY{p}{)}
\PY{n}{hc\PYZus{}comp\PYZus{}scale} \PY{o}{=} \PY{n}{HClust}\PY{p}{(}\PY{n}{distance\PYZus{}threshold}\PY{o}{=}\PY{l+m+mi}{0}\PY{p}{,}
                       \PY{n}{n\PYZus{}clusters}\PY{o}{=}\PY{k+kc}{None}\PY{p}{,}
                       \PY{n}{linkage}\PY{o}{=}\PY{l+s+s1}{\PYZsq{}}\PY{l+s+s1}{complete}\PY{l+s+s1}{\PYZsq{}}\PY{p}{)}\PY{o}{.}\PY{n}{fit}\PY{p}{(}\PY{n}{X\PYZus{}scale}\PY{p}{)}
\PY{n}{linkage\PYZus{}comp\PYZus{}scale} \PY{o}{=} \PY{n}{compute\PYZus{}linkage}\PY{p}{(}\PY{n}{hc\PYZus{}comp\PYZus{}scale}\PY{p}{)}
\PY{n}{fig}\PY{p}{,} \PY{n}{ax} \PY{o}{=} \PY{n}{plt}\PY{o}{.}\PY{n}{subplots}\PY{p}{(}\PY{l+m+mi}{1}\PY{p}{,} \PY{l+m+mi}{1}\PY{p}{,} \PY{n}{figsize}\PY{o}{=}\PY{p}{(}\PY{l+m+mi}{8}\PY{p}{,} \PY{l+m+mi}{8}\PY{p}{)}\PY{p}{)}
\PY{n}{dendrogram}\PY{p}{(}\PY{n}{linkage\PYZus{}comp\PYZus{}scale}\PY{p}{,} \PY{n}{ax}\PY{o}{=}\PY{n}{ax}\PY{p}{,} \PY{o}{*}\PY{o}{*}\PY{n}{cargs}\PY{p}{)}
\PY{n}{ax}\PY{o}{.}\PY{n}{set\PYZus{}title}\PY{p}{(}\PY{l+s+s2}{\PYZdq{}}\PY{l+s+s2}{Hierarchical Clustering with Scaled Features}\PY{l+s+s2}{\PYZdq{}}\PY{p}{)}\PY{p}{;}
\end{Verbatim}
\end{tcolorbox}

    \begin{center}
    \adjustimage{max size={0.9\linewidth}{0.9\paperheight}}{artifacts/latex/Ch12-unsup-lab_files/artifacts/latex/Ch12-unsup-lab_91_0.png}
    \end{center}
    { \hspace*{\fill} \\}
    
    Correlation-based distances between observations can be used for
clustering. The correlation between two observations measures the
similarity of their feature values. \{Suppose each observation has \(p\)
features, each a single numerical value. We measure the similarity of
two such observations by computing the correlation of these \(p\) pairs
of numbers.\} With \(n\) observations, the \(n\times n\) correlation
matrix can then be used as a similarity (or affinity) matrix, i.e.~so
that one minus the correlation matrix is the dissimilarity matrix used
for clustering.

Note that using correlation only makes sense for data with at least
three features since the absolute correlation between any two
observations with measurements on two features is always one. Hence, we
will cluster a three-dimensional data set.

    \begin{tcolorbox}[breakable, size=fbox, boxrule=1pt, pad at break*=1mm,colback=cellbackground, colframe=cellborder]
\prompt{In}{incolor}{46}{\boxspacing}
\begin{Verbatim}[commandchars=\\\{\}]
\PY{n}{X} \PY{o}{=} \PY{n}{np}\PY{o}{.}\PY{n}{random}\PY{o}{.}\PY{n}{standard\PYZus{}normal}\PY{p}{(}\PY{p}{(}\PY{l+m+mi}{30}\PY{p}{,} \PY{l+m+mi}{3}\PY{p}{)}\PY{p}{)}
\PY{n}{corD} \PY{o}{=} \PY{l+m+mi}{1} \PY{o}{\PYZhy{}} \PY{n}{np}\PY{o}{.}\PY{n}{corrcoef}\PY{p}{(}\PY{n}{X}\PY{p}{)}
\PY{n}{hc\PYZus{}cor} \PY{o}{=} \PY{n}{HClust}\PY{p}{(}\PY{n}{linkage}\PY{o}{=}\PY{l+s+s1}{\PYZsq{}}\PY{l+s+s1}{complete}\PY{l+s+s1}{\PYZsq{}}\PY{p}{,}
                \PY{n}{distance\PYZus{}threshold}\PY{o}{=}\PY{l+m+mi}{0}\PY{p}{,}
                \PY{n}{n\PYZus{}clusters}\PY{o}{=}\PY{k+kc}{None}\PY{p}{,}
                \PY{n}{metric}\PY{o}{=}\PY{l+s+s1}{\PYZsq{}}\PY{l+s+s1}{precomputed}\PY{l+s+s1}{\PYZsq{}}\PY{p}{)}
\PY{n}{hc\PYZus{}cor}\PY{o}{.}\PY{n}{fit}\PY{p}{(}\PY{n}{corD}\PY{p}{)}
\PY{n}{linkage\PYZus{}cor} \PY{o}{=} \PY{n}{compute\PYZus{}linkage}\PY{p}{(}\PY{n}{hc\PYZus{}cor}\PY{p}{)}
\PY{n}{fig}\PY{p}{,} \PY{n}{ax} \PY{o}{=} \PY{n}{plt}\PY{o}{.}\PY{n}{subplots}\PY{p}{(}\PY{l+m+mi}{1}\PY{p}{,} \PY{l+m+mi}{1}\PY{p}{,} \PY{n}{figsize}\PY{o}{=}\PY{p}{(}\PY{l+m+mi}{8}\PY{p}{,} \PY{l+m+mi}{8}\PY{p}{)}\PY{p}{)}
\PY{n}{dendrogram}\PY{p}{(}\PY{n}{linkage\PYZus{}cor}\PY{p}{,} \PY{n}{ax}\PY{o}{=}\PY{n}{ax}\PY{p}{,} \PY{o}{*}\PY{o}{*}\PY{n}{cargs}\PY{p}{)}
\PY{n}{ax}\PY{o}{.}\PY{n}{set\PYZus{}title}\PY{p}{(}\PY{l+s+s2}{\PYZdq{}}\PY{l+s+s2}{Complete Linkage with Correlation\PYZhy{}Based Dissimilarity}\PY{l+s+s2}{\PYZdq{}}\PY{p}{)}\PY{p}{;}
\end{Verbatim}
\end{tcolorbox}

    \begin{center}
    \adjustimage{max size={0.9\linewidth}{0.9\paperheight}}{artifacts/latex/Ch12-unsup-lab_files/artifacts/latex/Ch12-unsup-lab_93_0.png}
    \end{center}
    { \hspace*{\fill} \\}
    
    \hypertarget{nci60-data-example}{%
\subsection{NCI60 Data Example}\label{nci60-data-example}}

Unsupervised techniques are often used in the analysis of genomic data.
In particular, PCA and hierarchical clustering are popular tools. We
illustrate these techniques on the \texttt{NCI60} cancer cell line
microarray data, which consists of 6830 gene expression measurements on
64 cancer cell lines.

    \begin{tcolorbox}[breakable, size=fbox, boxrule=1pt, pad at break*=1mm,colback=cellbackground, colframe=cellborder]
\prompt{In}{incolor}{47}{\boxspacing}
\begin{Verbatim}[commandchars=\\\{\}]
\PY{n}{NCI60} \PY{o}{=} \PY{n}{load\PYZus{}data}\PY{p}{(}\PY{l+s+s1}{\PYZsq{}}\PY{l+s+s1}{NCI60}\PY{l+s+s1}{\PYZsq{}}\PY{p}{)}
\PY{n}{nci\PYZus{}labs} \PY{o}{=} \PY{n}{NCI60}\PY{p}{[}\PY{l+s+s1}{\PYZsq{}}\PY{l+s+s1}{labels}\PY{l+s+s1}{\PYZsq{}}\PY{p}{]}
\PY{n}{nci\PYZus{}data} \PY{o}{=} \PY{n}{NCI60}\PY{p}{[}\PY{l+s+s1}{\PYZsq{}}\PY{l+s+s1}{data}\PY{l+s+s1}{\PYZsq{}}\PY{p}{]}
\end{Verbatim}
\end{tcolorbox}

    Each cell line is labeled with a cancer type. We do not make use of the
cancer types in performing PCA and clustering, as these are unsupervised
techniques. But after performing PCA and clustering, we will check to
see the extent to which these cancer types agree with the results of
these unsupervised techniques.

The data has 64 rows and 6830 columns.

    \begin{tcolorbox}[breakable, size=fbox, boxrule=1pt, pad at break*=1mm,colback=cellbackground, colframe=cellborder]
\prompt{In}{incolor}{48}{\boxspacing}
\begin{Verbatim}[commandchars=\\\{\}]
\PY{n}{nci\PYZus{}data}\PY{o}{.}\PY{n}{shape}
\end{Verbatim}
\end{tcolorbox}

            \begin{tcolorbox}[breakable, size=fbox, boxrule=.5pt, pad at break*=1mm, opacityfill=0]
\prompt{Out}{outcolor}{48}{\boxspacing}
\begin{Verbatim}[commandchars=\\\{\}]
(64, 6830)
\end{Verbatim}
\end{tcolorbox}
        
    We begin by examining the cancer types for the cell lines.

    \begin{tcolorbox}[breakable, size=fbox, boxrule=1pt, pad at break*=1mm,colback=cellbackground, colframe=cellborder]
\prompt{In}{incolor}{49}{\boxspacing}
\begin{Verbatim}[commandchars=\\\{\}]
\PY{n}{nci\PYZus{}labs}\PY{o}{.}\PY{n}{value\PYZus{}counts}\PY{p}{(}\PY{p}{)}
\end{Verbatim}
\end{tcolorbox}

            \begin{tcolorbox}[breakable, size=fbox, boxrule=.5pt, pad at break*=1mm, opacityfill=0]
\prompt{Out}{outcolor}{49}{\boxspacing}
\begin{Verbatim}[commandchars=\\\{\}]
label
NSCLC          9
RENAL          9
MELANOMA       8
BREAST         7
COLON          7
LEUKEMIA       6
OVARIAN        6
CNS            5
PROSTATE       2
K562A-repro    1
K562B-repro    1
MCF7A-repro    1
MCF7D-repro    1
UNKNOWN        1
Name: count, dtype: int64
\end{Verbatim}
\end{tcolorbox}
        
    \hypertarget{pca-on-the-nci60-data}{%
\subsubsection{PCA on the NCI60 Data}\label{pca-on-the-nci60-data}}

We first perform PCA on the data after scaling the variables (genes) to
have standard deviation one, although here one could reasonably argue
that it is better not to scale the genes as they are measured in the
same units.

    \begin{tcolorbox}[breakable, size=fbox, boxrule=1pt, pad at break*=1mm,colback=cellbackground, colframe=cellborder]
\prompt{In}{incolor}{50}{\boxspacing}
\begin{Verbatim}[commandchars=\\\{\}]
\PY{n}{scaler} \PY{o}{=} \PY{n}{StandardScaler}\PY{p}{(}\PY{p}{)}
\PY{n}{nci\PYZus{}scaled} \PY{o}{=} \PY{n}{scaler}\PY{o}{.}\PY{n}{fit\PYZus{}transform}\PY{p}{(}\PY{n}{nci\PYZus{}data}\PY{p}{)}
\PY{n}{nci\PYZus{}pca} \PY{o}{=} \PY{n}{PCA}\PY{p}{(}\PY{p}{)}
\PY{n}{nci\PYZus{}scores} \PY{o}{=} \PY{n}{nci\PYZus{}pca}\PY{o}{.}\PY{n}{fit\PYZus{}transform}\PY{p}{(}\PY{n}{nci\PYZus{}scaled}\PY{p}{)}
\end{Verbatim}
\end{tcolorbox}

    We now plot the first few principal component score vectors, in order to
visualize the data. The observations (cell lines) corresponding to a
given cancer type will be plotted in the same color, so that we can see
to what extent the observations within a cancer type are similar to each
other.

    \begin{tcolorbox}[breakable, size=fbox, boxrule=1pt, pad at break*=1mm,colback=cellbackground, colframe=cellborder]
\prompt{In}{incolor}{51}{\boxspacing}
\begin{Verbatim}[commandchars=\\\{\}]
\PY{n}{cancer\PYZus{}types} \PY{o}{=} \PY{n+nb}{list}\PY{p}{(}\PY{n}{np}\PY{o}{.}\PY{n}{unique}\PY{p}{(}\PY{n}{nci\PYZus{}labs}\PY{p}{)}\PY{p}{)}
\PY{n}{nci\PYZus{}groups} \PY{o}{=} \PY{n}{np}\PY{o}{.}\PY{n}{array}\PY{p}{(}\PY{p}{[}\PY{n}{cancer\PYZus{}types}\PY{o}{.}\PY{n}{index}\PY{p}{(}\PY{n}{lab}\PY{p}{)}
                       \PY{k}{for} \PY{n}{lab} \PY{o+ow}{in} \PY{n}{nci\PYZus{}labs}\PY{o}{.}\PY{n}{values}\PY{p}{]}\PY{p}{)}
\PY{n}{fig}\PY{p}{,} \PY{n}{axes} \PY{o}{=} \PY{n}{plt}\PY{o}{.}\PY{n}{subplots}\PY{p}{(}\PY{l+m+mi}{1}\PY{p}{,} \PY{l+m+mi}{2}\PY{p}{,} \PY{n}{figsize}\PY{o}{=}\PY{p}{(}\PY{l+m+mi}{15}\PY{p}{,}\PY{l+m+mi}{6}\PY{p}{)}\PY{p}{)}
\PY{n}{ax} \PY{o}{=} \PY{n}{axes}\PY{p}{[}\PY{l+m+mi}{0}\PY{p}{]}
\PY{n}{ax}\PY{o}{.}\PY{n}{scatter}\PY{p}{(}\PY{n}{nci\PYZus{}scores}\PY{p}{[}\PY{p}{:}\PY{p}{,}\PY{l+m+mi}{0}\PY{p}{]}\PY{p}{,}
           \PY{n}{nci\PYZus{}scores}\PY{p}{[}\PY{p}{:}\PY{p}{,}\PY{l+m+mi}{1}\PY{p}{]}\PY{p}{,}
           \PY{n}{c}\PY{o}{=}\PY{n}{nci\PYZus{}groups}\PY{p}{,}
           \PY{n}{marker}\PY{o}{=}\PY{l+s+s1}{\PYZsq{}}\PY{l+s+s1}{o}\PY{l+s+s1}{\PYZsq{}}\PY{p}{,}
           \PY{n}{s}\PY{o}{=}\PY{l+m+mi}{50}\PY{p}{)}
\PY{n}{ax}\PY{o}{.}\PY{n}{set\PYZus{}xlabel}\PY{p}{(}\PY{l+s+s1}{\PYZsq{}}\PY{l+s+s1}{PC1}\PY{l+s+s1}{\PYZsq{}}\PY{p}{)}\PY{p}{;} \PY{n}{ax}\PY{o}{.}\PY{n}{set\PYZus{}ylabel}\PY{p}{(}\PY{l+s+s1}{\PYZsq{}}\PY{l+s+s1}{PC2}\PY{l+s+s1}{\PYZsq{}}\PY{p}{)}
\PY{n}{ax} \PY{o}{=} \PY{n}{axes}\PY{p}{[}\PY{l+m+mi}{1}\PY{p}{]}
\PY{n}{ax}\PY{o}{.}\PY{n}{scatter}\PY{p}{(}\PY{n}{nci\PYZus{}scores}\PY{p}{[}\PY{p}{:}\PY{p}{,}\PY{l+m+mi}{0}\PY{p}{]}\PY{p}{,}
           \PY{n}{nci\PYZus{}scores}\PY{p}{[}\PY{p}{:}\PY{p}{,}\PY{l+m+mi}{2}\PY{p}{]}\PY{p}{,}
           \PY{n}{c}\PY{o}{=}\PY{n}{nci\PYZus{}groups}\PY{p}{,}
           \PY{n}{marker}\PY{o}{=}\PY{l+s+s1}{\PYZsq{}}\PY{l+s+s1}{o}\PY{l+s+s1}{\PYZsq{}}\PY{p}{,}
           \PY{n}{s}\PY{o}{=}\PY{l+m+mi}{50}\PY{p}{)}
\PY{n}{ax}\PY{o}{.}\PY{n}{set\PYZus{}xlabel}\PY{p}{(}\PY{l+s+s1}{\PYZsq{}}\PY{l+s+s1}{PC1}\PY{l+s+s1}{\PYZsq{}}\PY{p}{)}\PY{p}{;} \PY{n}{ax}\PY{o}{.}\PY{n}{set\PYZus{}ylabel}\PY{p}{(}\PY{l+s+s1}{\PYZsq{}}\PY{l+s+s1}{PC3}\PY{l+s+s1}{\PYZsq{}}\PY{p}{)}\PY{p}{;}
\end{Verbatim}
\end{tcolorbox}

    \begin{center}
    \adjustimage{max size={0.9\linewidth}{0.9\paperheight}}{artifacts/latex/Ch12-unsup-lab_files/artifacts/latex/Ch12-unsup-lab_103_0.png}
    \end{center}
    { \hspace*{\fill} \\}
    
    On the whole, cell lines corresponding to a single cancer type do tend
to have similar values on the first few principal component score
vectors. This indicates that cell lines from the same cancer type tend
to have pretty similar gene expression levels.

    We can also plot the percent variance explained by the principal
components as well as the cumulative percent variance explained. This is
similar to the plots we made earlier for the \texttt{USArrests} data.

    \begin{tcolorbox}[breakable, size=fbox, boxrule=1pt, pad at break*=1mm,colback=cellbackground, colframe=cellborder]
\prompt{In}{incolor}{52}{\boxspacing}
\begin{Verbatim}[commandchars=\\\{\}]
\PY{n}{fig}\PY{p}{,} \PY{n}{axes} \PY{o}{=} \PY{n}{plt}\PY{o}{.}\PY{n}{subplots}\PY{p}{(}\PY{l+m+mi}{1}\PY{p}{,} \PY{l+m+mi}{2}\PY{p}{,} \PY{n}{figsize}\PY{o}{=}\PY{p}{(}\PY{l+m+mi}{15}\PY{p}{,}\PY{l+m+mi}{6}\PY{p}{)}\PY{p}{)}
\PY{n}{ax} \PY{o}{=} \PY{n}{axes}\PY{p}{[}\PY{l+m+mi}{0}\PY{p}{]}
\PY{n}{ticks} \PY{o}{=} \PY{n}{np}\PY{o}{.}\PY{n}{arange}\PY{p}{(}\PY{n}{nci\PYZus{}pca}\PY{o}{.}\PY{n}{n\PYZus{}components\PYZus{}}\PY{p}{)}\PY{o}{+}\PY{l+m+mi}{1}
\PY{n}{ax}\PY{o}{.}\PY{n}{plot}\PY{p}{(}\PY{n}{ticks}\PY{p}{,}
        \PY{n}{nci\PYZus{}pca}\PY{o}{.}\PY{n}{explained\PYZus{}variance\PYZus{}ratio\PYZus{}}\PY{p}{,}
        \PY{n}{marker}\PY{o}{=}\PY{l+s+s1}{\PYZsq{}}\PY{l+s+s1}{o}\PY{l+s+s1}{\PYZsq{}}\PY{p}{)}
\PY{n}{ax}\PY{o}{.}\PY{n}{set\PYZus{}xlabel}\PY{p}{(}\PY{l+s+s1}{\PYZsq{}}\PY{l+s+s1}{Principal Component}\PY{l+s+s1}{\PYZsq{}}\PY{p}{)}\PY{p}{;}
\PY{n}{ax}\PY{o}{.}\PY{n}{set\PYZus{}ylabel}\PY{p}{(}\PY{l+s+s1}{\PYZsq{}}\PY{l+s+s1}{PVE}\PY{l+s+s1}{\PYZsq{}}\PY{p}{)}
\PY{n}{ax} \PY{o}{=} \PY{n}{axes}\PY{p}{[}\PY{l+m+mi}{1}\PY{p}{]}
\PY{n}{ax}\PY{o}{.}\PY{n}{plot}\PY{p}{(}\PY{n}{ticks}\PY{p}{,}
        \PY{n}{nci\PYZus{}pca}\PY{o}{.}\PY{n}{explained\PYZus{}variance\PYZus{}ratio\PYZus{}}\PY{o}{.}\PY{n}{cumsum}\PY{p}{(}\PY{p}{)}\PY{p}{,}
        \PY{n}{marker}\PY{o}{=}\PY{l+s+s1}{\PYZsq{}}\PY{l+s+s1}{o}\PY{l+s+s1}{\PYZsq{}}\PY{p}{)}\PY{p}{;}
\PY{n}{ax}\PY{o}{.}\PY{n}{set\PYZus{}xlabel}\PY{p}{(}\PY{l+s+s1}{\PYZsq{}}\PY{l+s+s1}{Principal Component}\PY{l+s+s1}{\PYZsq{}}\PY{p}{)}
\PY{n}{ax}\PY{o}{.}\PY{n}{set\PYZus{}ylabel}\PY{p}{(}\PY{l+s+s1}{\PYZsq{}}\PY{l+s+s1}{Cumulative PVE}\PY{l+s+s1}{\PYZsq{}}\PY{p}{)}\PY{p}{;}
\end{Verbatim}
\end{tcolorbox}

    \begin{center}
    \adjustimage{max size={0.9\linewidth}{0.9\paperheight}}{artifacts/latex/Ch12-unsup-lab_files/artifacts/latex/Ch12-unsup-lab_106_0.png}
    \end{center}
    { \hspace*{\fill} \\}
    
    We see that together, the first seven principal components explain
around 40\% of the variance in the data. This is not a huge amount of
the variance. However, looking at the scree plot, we see that while each
of the first seven principal components explain a substantial amount of
variance, there is a marked decrease in the variance explained by
further principal components. That is, there is an \emph{elbow} in the
plot after approximately the seventh principal component. This suggests
that there may be little benefit to examining more than seven or so
principal components (though even examining seven principal components
may be difficult).

    \hypertarget{clustering-the-observations-of-the-nci60-data}{%
\subsubsection{Clustering the Observations of the NCI60
Data}\label{clustering-the-observations-of-the-nci60-data}}

We now perform hierarchical clustering of the cell lines in the
\texttt{NCI60} data using complete, single, and average linkage. Once
again, the goal is to find out whether or not the observations cluster
into distinct types of cancer. Euclidean distance is used as the
dissimilarity measure. We first write a short function to produce the
three dendrograms.

    \begin{tcolorbox}[breakable, size=fbox, boxrule=1pt, pad at break*=1mm,colback=cellbackground, colframe=cellborder]
\prompt{In}{incolor}{53}{\boxspacing}
\begin{Verbatim}[commandchars=\\\{\}]
\PY{k}{def}\PY{+w}{ }\PY{n+nf}{plot\PYZus{}nci}\PY{p}{(}\PY{n}{linkage}\PY{p}{,} \PY{n}{ax}\PY{p}{,} \PY{n}{cut}\PY{o}{=}\PY{o}{\PYZhy{}}\PY{n}{np}\PY{o}{.}\PY{n}{inf}\PY{p}{)}\PY{p}{:}
    \PY{n}{cargs} \PY{o}{=} \PY{p}{\PYZob{}}\PY{l+s+s1}{\PYZsq{}}\PY{l+s+s1}{above\PYZus{}threshold\PYZus{}color}\PY{l+s+s1}{\PYZsq{}}\PY{p}{:}\PY{l+s+s1}{\PYZsq{}}\PY{l+s+s1}{black}\PY{l+s+s1}{\PYZsq{}}\PY{p}{,}
             \PY{l+s+s1}{\PYZsq{}}\PY{l+s+s1}{color\PYZus{}threshold}\PY{l+s+s1}{\PYZsq{}}\PY{p}{:}\PY{n}{cut}\PY{p}{\PYZcb{}}
    \PY{n}{hc} \PY{o}{=} \PY{n}{HClust}\PY{p}{(}\PY{n}{n\PYZus{}clusters}\PY{o}{=}\PY{k+kc}{None}\PY{p}{,}
                \PY{n}{distance\PYZus{}threshold}\PY{o}{=}\PY{l+m+mi}{0}\PY{p}{,}
                \PY{n}{linkage}\PY{o}{=}\PY{n}{linkage}\PY{o}{.}\PY{n}{lower}\PY{p}{(}\PY{p}{)}\PY{p}{)}\PY{o}{.}\PY{n}{fit}\PY{p}{(}\PY{n}{nci\PYZus{}scaled}\PY{p}{)}
    \PY{n}{linkage\PYZus{}} \PY{o}{=} \PY{n}{compute\PYZus{}linkage}\PY{p}{(}\PY{n}{hc}\PY{p}{)}
    \PY{n}{dendrogram}\PY{p}{(}\PY{n}{linkage\PYZus{}}\PY{p}{,}
               \PY{n}{ax}\PY{o}{=}\PY{n}{ax}\PY{p}{,}
               \PY{n}{labels}\PY{o}{=}\PY{n}{np}\PY{o}{.}\PY{n}{asarray}\PY{p}{(}\PY{n}{nci\PYZus{}labs}\PY{p}{)}\PY{p}{,}
               \PY{n}{leaf\PYZus{}font\PYZus{}size}\PY{o}{=}\PY{l+m+mi}{10}\PY{p}{,}
               \PY{o}{*}\PY{o}{*}\PY{n}{cargs}\PY{p}{)}
    \PY{n}{ax}\PY{o}{.}\PY{n}{set\PYZus{}title}\PY{p}{(}\PY{l+s+s1}{\PYZsq{}}\PY{l+s+si}{\PYZpc{}s}\PY{l+s+s1}{ Linkage}\PY{l+s+s1}{\PYZsq{}} \PY{o}{\PYZpc{}} \PY{n}{linkage}\PY{p}{)}
    \PY{k}{return} \PY{n}{hc}
\end{Verbatim}
\end{tcolorbox}

    Let's plot our results.

    \begin{tcolorbox}[breakable, size=fbox, boxrule=1pt, pad at break*=1mm,colback=cellbackground, colframe=cellborder]
\prompt{In}{incolor}{54}{\boxspacing}
\begin{Verbatim}[commandchars=\\\{\}]
\PY{n}{fig}\PY{p}{,} \PY{n}{axes} \PY{o}{=} \PY{n}{plt}\PY{o}{.}\PY{n}{subplots}\PY{p}{(}\PY{l+m+mi}{3}\PY{p}{,} \PY{l+m+mi}{1}\PY{p}{,} \PY{n}{figsize}\PY{o}{=}\PY{p}{(}\PY{l+m+mi}{15}\PY{p}{,}\PY{l+m+mi}{30}\PY{p}{)}\PY{p}{)}      
\PY{n}{ax} \PY{o}{=} \PY{n}{axes}\PY{p}{[}\PY{l+m+mi}{0}\PY{p}{]}\PY{p}{;} \PY{n}{hc\PYZus{}comp} \PY{o}{=} \PY{n}{plot\PYZus{}nci}\PY{p}{(}\PY{l+s+s1}{\PYZsq{}}\PY{l+s+s1}{Complete}\PY{l+s+s1}{\PYZsq{}}\PY{p}{,} \PY{n}{ax}\PY{p}{)}
\PY{n}{ax} \PY{o}{=} \PY{n}{axes}\PY{p}{[}\PY{l+m+mi}{1}\PY{p}{]}\PY{p}{;} \PY{n}{hc\PYZus{}avg} \PY{o}{=} \PY{n}{plot\PYZus{}nci}\PY{p}{(}\PY{l+s+s1}{\PYZsq{}}\PY{l+s+s1}{Average}\PY{l+s+s1}{\PYZsq{}}\PY{p}{,} \PY{n}{ax}\PY{p}{)}
\PY{n}{ax} \PY{o}{=} \PY{n}{axes}\PY{p}{[}\PY{l+m+mi}{2}\PY{p}{]}\PY{p}{;} \PY{n}{hc\PYZus{}sing} \PY{o}{=} \PY{n}{plot\PYZus{}nci}\PY{p}{(}\PY{l+s+s1}{\PYZsq{}}\PY{l+s+s1}{Single}\PY{l+s+s1}{\PYZsq{}}\PY{p}{,} \PY{n}{ax}\PY{p}{)}
\end{Verbatim}
\end{tcolorbox}

    \begin{center}
    \adjustimage{max size={0.9\linewidth}{0.9\paperheight}}{artifacts/latex/Ch12-unsup-lab_files/artifacts/latex/Ch12-unsup-lab_111_0.png}
    \end{center}
    { \hspace*{\fill} \\}
    
    We see that the choice of linkage certainly does affect the results
obtained. Typically, single linkage will tend to yield \emph{trailing}
clusters: very large clusters onto which individual observations attach
one-by-one. On the other hand, complete and average linkage tend to
yield more balanced, attractive clusters. For this reason, complete and
average linkage are generally preferred to single linkage. Clearly cell
lines within a single cancer type do tend to cluster together, although
the clustering is not perfect. We will use complete linkage hierarchical
clustering for the analysis that follows.

We can cut the dendrogram at the height that will yield a particular
number of clusters, say four:

    \begin{tcolorbox}[breakable, size=fbox, boxrule=1pt, pad at break*=1mm,colback=cellbackground, colframe=cellborder]
\prompt{In}{incolor}{55}{\boxspacing}
\begin{Verbatim}[commandchars=\\\{\}]
\PY{n}{linkage\PYZus{}comp} \PY{o}{=} \PY{n}{compute\PYZus{}linkage}\PY{p}{(}\PY{n}{hc\PYZus{}comp}\PY{p}{)}
\PY{n}{comp\PYZus{}cut} \PY{o}{=} \PY{n}{cut\PYZus{}tree}\PY{p}{(}\PY{n}{linkage\PYZus{}comp}\PY{p}{,} \PY{n}{n\PYZus{}clusters}\PY{o}{=}\PY{l+m+mi}{4}\PY{p}{)}\PY{o}{.}\PY{n}{reshape}\PY{p}{(}\PY{o}{\PYZhy{}}\PY{l+m+mi}{1}\PY{p}{)}
\PY{n}{pd}\PY{o}{.}\PY{n}{crosstab}\PY{p}{(}\PY{n}{nci\PYZus{}labs}\PY{p}{[}\PY{l+s+s1}{\PYZsq{}}\PY{l+s+s1}{label}\PY{l+s+s1}{\PYZsq{}}\PY{p}{]}\PY{p}{,}
            \PY{n}{pd}\PY{o}{.}\PY{n}{Series}\PY{p}{(}\PY{n}{comp\PYZus{}cut}\PY{o}{.}\PY{n}{reshape}\PY{p}{(}\PY{o}{\PYZhy{}}\PY{l+m+mi}{1}\PY{p}{)}\PY{p}{,} \PY{n}{name}\PY{o}{=}\PY{l+s+s1}{\PYZsq{}}\PY{l+s+s1}{Complete}\PY{l+s+s1}{\PYZsq{}}\PY{p}{)}\PY{p}{)}
\end{Verbatim}
\end{tcolorbox}

            \begin{tcolorbox}[breakable, size=fbox, boxrule=.5pt, pad at break*=1mm, opacityfill=0]
\prompt{Out}{outcolor}{55}{\boxspacing}
\begin{Verbatim}[commandchars=\\\{\}]
Complete     0  1  2  3
label
BREAST       2  3  0  2
CNS          3  2  0  0
COLON        2  0  0  5
K562A-repro  0  0  1  0
K562B-repro  0  0  1  0
LEUKEMIA     0  0  6  0
MCF7A-repro  0  0  0  1
MCF7D-repro  0  0  0  1
MELANOMA     8  0  0  0
NSCLC        8  1  0  0
OVARIAN      6  0  0  0
PROSTATE     2  0  0  0
RENAL        8  1  0  0
UNKNOWN      1  0  0  0
\end{Verbatim}
\end{tcolorbox}
        
    There are some clear patterns. All the leukemia cell lines fall in one
cluster, while the breast cancer cell lines are spread out over three
different clusters.

We can plot a cut on the dendrogram that produces these four clusters:

    \begin{tcolorbox}[breakable, size=fbox, boxrule=1pt, pad at break*=1mm,colback=cellbackground, colframe=cellborder]
\prompt{In}{incolor}{56}{\boxspacing}
\begin{Verbatim}[commandchars=\\\{\}]
\PY{n}{fig}\PY{p}{,} \PY{n}{ax} \PY{o}{=} \PY{n}{plt}\PY{o}{.}\PY{n}{subplots}\PY{p}{(}\PY{n}{figsize}\PY{o}{=}\PY{p}{(}\PY{l+m+mi}{10}\PY{p}{,}\PY{l+m+mi}{10}\PY{p}{)}\PY{p}{)}
\PY{n}{plot\PYZus{}nci}\PY{p}{(}\PY{l+s+s1}{\PYZsq{}}\PY{l+s+s1}{Complete}\PY{l+s+s1}{\PYZsq{}}\PY{p}{,} \PY{n}{ax}\PY{p}{,} \PY{n}{cut}\PY{o}{=}\PY{l+m+mi}{140}\PY{p}{)}
\PY{n}{ax}\PY{o}{.}\PY{n}{axhline}\PY{p}{(}\PY{l+m+mi}{140}\PY{p}{,} \PY{n}{c}\PY{o}{=}\PY{l+s+s1}{\PYZsq{}}\PY{l+s+s1}{r}\PY{l+s+s1}{\PYZsq{}}\PY{p}{,} \PY{n}{linewidth}\PY{o}{=}\PY{l+m+mi}{4}\PY{p}{)}\PY{p}{;}
\end{Verbatim}
\end{tcolorbox}

    \begin{center}
    \adjustimage{max size={0.9\linewidth}{0.9\paperheight}}{artifacts/latex/Ch12-unsup-lab_files/artifacts/latex/Ch12-unsup-lab_115_0.png}
    \end{center}
    { \hspace*{\fill} \\}
    
    The \texttt{axhline()} function draws a horizontal line on top of any
existing set of axes. The argument \texttt{140} plots a horizontal line
at height 140 on the dendrogram; this is a height that results in four
distinct clusters. It is easy to verify that the resulting clusters are
the same as the ones we obtained in \texttt{comp\_cut}.

We claimed earlier in Section 12.4.2 that \(K\)-means clustering and
hierarchical clustering with the dendrogram cut to obtain the same
number of clusters can yield very different results. How do these
\texttt{NCI60} hierarchical clustering results compare to what we get if
we perform \(K\)-means clustering with \(K=4\)?

    \begin{tcolorbox}[breakable, size=fbox, boxrule=1pt, pad at break*=1mm,colback=cellbackground, colframe=cellborder]
\prompt{In}{incolor}{57}{\boxspacing}
\begin{Verbatim}[commandchars=\\\{\}]
\PY{n}{nci\PYZus{}kmeans} \PY{o}{=} \PY{n}{KMeans}\PY{p}{(}\PY{n}{n\PYZus{}clusters}\PY{o}{=}\PY{l+m+mi}{4}\PY{p}{,} 
                    \PY{n}{random\PYZus{}state}\PY{o}{=}\PY{l+m+mi}{0}\PY{p}{,}
                    \PY{n}{n\PYZus{}init}\PY{o}{=}\PY{l+m+mi}{20}\PY{p}{)}\PY{o}{.}\PY{n}{fit}\PY{p}{(}\PY{n}{nci\PYZus{}scaled}\PY{p}{)}
\PY{n}{pd}\PY{o}{.}\PY{n}{crosstab}\PY{p}{(}\PY{n}{pd}\PY{o}{.}\PY{n}{Series}\PY{p}{(}\PY{n}{comp\PYZus{}cut}\PY{p}{,} \PY{n}{name}\PY{o}{=}\PY{l+s+s1}{\PYZsq{}}\PY{l+s+s1}{HClust}\PY{l+s+s1}{\PYZsq{}}\PY{p}{)}\PY{p}{,}
            \PY{n}{pd}\PY{o}{.}\PY{n}{Series}\PY{p}{(}\PY{n}{nci\PYZus{}kmeans}\PY{o}{.}\PY{n}{labels\PYZus{}}\PY{p}{,} \PY{n}{name}\PY{o}{=}\PY{l+s+s1}{\PYZsq{}}\PY{l+s+s1}{K\PYZhy{}means}\PY{l+s+s1}{\PYZsq{}}\PY{p}{)}\PY{p}{)}
\end{Verbatim}
\end{tcolorbox}

            \begin{tcolorbox}[breakable, size=fbox, boxrule=.5pt, pad at break*=1mm, opacityfill=0]
\prompt{Out}{outcolor}{57}{\boxspacing}
\begin{Verbatim}[commandchars=\\\{\}]
K-means  0   1   2  3
HClust
0        1  20  10  9
1        0   7   0  0
2        8   0   0  0
3        0   0   9  0
\end{Verbatim}
\end{tcolorbox}
        
    We see that the four clusters obtained using hierarchical clustering and
\(K\)-means clustering are somewhat different. First we note that the
labels in the two clusterings are arbitrary. That is, swapping the
identifier of the cluster does not change the clustering. We see here
Cluster 3 in \(K\)-means clustering is identical to cluster 2 in
hierarchical clustering. However, the other clusters differ: for
instance, cluster 0 in \(K\)-means clustering contains a portion of the
observations assigned to cluster 0 by hierarchical clustering, as well
as all of the observations assigned to cluster 1 by hierarchical
clustering.

Rather than performing hierarchical clustering on the entire data
matrix, we can also perform hierarchical clustering on the first few
principal component score vectors, regarding these first few components
as a less noisy version of the data.

    \begin{tcolorbox}[breakable, size=fbox, boxrule=1pt, pad at break*=1mm,colback=cellbackground, colframe=cellborder]
\prompt{In}{incolor}{58}{\boxspacing}
\begin{Verbatim}[commandchars=\\\{\}]
\PY{n}{hc\PYZus{}pca} \PY{o}{=} \PY{n}{HClust}\PY{p}{(}\PY{n}{n\PYZus{}clusters}\PY{o}{=}\PY{k+kc}{None}\PY{p}{,}
                \PY{n}{distance\PYZus{}threshold}\PY{o}{=}\PY{l+m+mi}{0}\PY{p}{,}
                \PY{n}{linkage}\PY{o}{=}\PY{l+s+s1}{\PYZsq{}}\PY{l+s+s1}{complete}\PY{l+s+s1}{\PYZsq{}}
                \PY{p}{)}\PY{o}{.}\PY{n}{fit}\PY{p}{(}\PY{n}{nci\PYZus{}scores}\PY{p}{[}\PY{p}{:}\PY{p}{,}\PY{p}{:}\PY{l+m+mi}{5}\PY{p}{]}\PY{p}{)}
\PY{n}{linkage\PYZus{}pca} \PY{o}{=} \PY{n}{compute\PYZus{}linkage}\PY{p}{(}\PY{n}{hc\PYZus{}pca}\PY{p}{)}
\PY{n}{fig}\PY{p}{,} \PY{n}{ax} \PY{o}{=} \PY{n}{plt}\PY{o}{.}\PY{n}{subplots}\PY{p}{(}\PY{n}{figsize}\PY{o}{=}\PY{p}{(}\PY{l+m+mi}{8}\PY{p}{,}\PY{l+m+mi}{8}\PY{p}{)}\PY{p}{)}
\PY{n}{dendrogram}\PY{p}{(}\PY{n}{linkage\PYZus{}pca}\PY{p}{,}
           \PY{n}{labels}\PY{o}{=}\PY{n}{np}\PY{o}{.}\PY{n}{asarray}\PY{p}{(}\PY{n}{nci\PYZus{}labs}\PY{p}{)}\PY{p}{,}
           \PY{n}{leaf\PYZus{}font\PYZus{}size}\PY{o}{=}\PY{l+m+mi}{10}\PY{p}{,}
           \PY{n}{ax}\PY{o}{=}\PY{n}{ax}\PY{p}{,}
           \PY{o}{*}\PY{o}{*}\PY{n}{cargs}\PY{p}{)}
\PY{n}{ax}\PY{o}{.}\PY{n}{set\PYZus{}title}\PY{p}{(}\PY{l+s+s2}{\PYZdq{}}\PY{l+s+s2}{Hier. Clust. on First Five Score Vectors}\PY{l+s+s2}{\PYZdq{}}\PY{p}{)}
\PY{n}{pca\PYZus{}labels} \PY{o}{=} \PY{n}{pd}\PY{o}{.}\PY{n}{Series}\PY{p}{(}\PY{n}{cut\PYZus{}tree}\PY{p}{(}\PY{n}{linkage\PYZus{}pca}\PY{p}{,}
                                \PY{n}{n\PYZus{}clusters}\PY{o}{=}\PY{l+m+mi}{4}\PY{p}{)}\PY{o}{.}\PY{n}{reshape}\PY{p}{(}\PY{o}{\PYZhy{}}\PY{l+m+mi}{1}\PY{p}{)}\PY{p}{,}
                       \PY{n}{name}\PY{o}{=}\PY{l+s+s1}{\PYZsq{}}\PY{l+s+s1}{Complete\PYZhy{}PCA}\PY{l+s+s1}{\PYZsq{}}\PY{p}{)}
\PY{n}{pd}\PY{o}{.}\PY{n}{crosstab}\PY{p}{(}\PY{n}{nci\PYZus{}labs}\PY{p}{[}\PY{l+s+s1}{\PYZsq{}}\PY{l+s+s1}{label}\PY{l+s+s1}{\PYZsq{}}\PY{p}{]}\PY{p}{,} \PY{n}{pca\PYZus{}labels}\PY{p}{)}
\end{Verbatim}
\end{tcolorbox}

            \begin{tcolorbox}[breakable, size=fbox, boxrule=.5pt, pad at break*=1mm, opacityfill=0]
\prompt{Out}{outcolor}{58}{\boxspacing}
\begin{Verbatim}[commandchars=\\\{\}]
Complete-PCA  0  1  2  3
label
BREAST        0  5  0  2
CNS           2  3  0  0
COLON         7  0  0  0
K562A-repro   0  0  1  0
K562B-repro   0  0  1  0
LEUKEMIA      2  0  4  0
MCF7A-repro   0  0  0  1
MCF7D-repro   0  0  0  1
MELANOMA      1  7  0  0
NSCLC         8  1  0  0
OVARIAN       5  1  0  0
PROSTATE      2  0  0  0
RENAL         7  2  0  0
UNKNOWN       0  1  0  0
\end{Verbatim}
\end{tcolorbox}
        
    \begin{center}
    \adjustimage{max size={0.9\linewidth}{0.9\paperheight}}{artifacts/latex/Ch12-unsup-lab_files/artifacts/latex/Ch12-unsup-lab_119_1.png}
    \end{center}
    { \hspace*{\fill} \\}
    
    


    % Add a bibliography block to the postdoc
    
    
    
\end{document}
